\documentclass[11pt]{article}

% Anonymous submission.
\usepackage[review]{acl}

\usepackage{times}
\usepackage{latexsym}
\usepackage[T1]{fontenc}
\usepackage[utf8]{inputenc}
\usepackage{microtype}
\usepackage{inconsolata}
\usepackage{graphicx}
\usepackage{placeins}
\usepackage{cuted}
\usepackage{capt-of}
\usepackage{booktabs}
\usepackage{multirow}
\usepackage{amsmath}
\usepackage{amssymb}
\usepackage{mathtools}
\usepackage{colortbl}
\usepackage{array}
\usepackage{url}
\usepackage{tcolorbox}
\tcbuselibrary{breakable}
\usepackage[ruled,vlined]{algorithm2e}

% Anonymous review shows only "Anonymous ACL submission", so the
% title box can be shorter. \titlebox is a fixed height.
\makeatletter
\ifacl@anonymize
  \setlength\titlebox{4.5cm}
\else
  \setlength\titlebox{7.6cm}
\fi
\makeatother

\newcommand{\ours}{\texttt{ExpG}}
\definecolor{gainbg}{RGB}{230,245,230}
\definecolor{darkgreen}{RGB}{0,110,60}
\definecolor{lightblue}{RGB}{102,167,208}
\definecolor{darkred}{RGB}{139,0,0}

\title{Constructing Experience-Driven Adaptive Guidance \\ for Continually Improving Tool Use in Agent Harnesses}

\author{Anonymous Authors}


\begin{document}


\maketitle


\begin{abstract}
The performance bottleneck of agents is increasingly shifting from model capability to the agent harness that governs execution. 
Within this harness, tools are the primary interface to external environments, yet available guidance rarely adapts across runtime conditions. 
We propose {\ours}, a mechanism that constructs experience-driven adaptive guidance capturing each tool's capability boundaries and best practices, so that tool use can improve as experience accumulates. 
{\ours} consists of three phases. 
(1) Experience acquisition analyzes tool invocation quality in historical trajectories through multi-aspect attribution. 
(2) Experience distillation filters unhelpful experiences, selects representative ones with an equivalence-class-based method, and summarizes them into guidance. 
(3) Experience reuse applies the guidance as contextual prompts or tool-schema constraints. 
Experiments show that {\ours} brings consistent improvements across tool selection, tool calling, and response generation, enabling smaller agents to outperform larger ones that do not use it. It achieves particularly strong gains in challenging settings, pointing toward continually improving tool use in agent harnesses. Our code is available in the \href{https://anonymous.4open.science/r/ExpG-F72B/README.md}{anonymous repository}.
\end{abstract}



\section{Introduction}


% By invoking external tools, agents can tap into domain-specific expertise, access continuously updated real-time data and interact with users in more flexible ways~\cite{1,3,5}. 
% Equipped with a rich ecosystem of external tools, agents overcome the limitations of their inherent knowledge and reasoning capabilities, thereby handling tasks with greater accuracy, reliability, and efficiency.
% Tool empowerment represents an important milestone in the evolution of large language models (LLMs) toward truly intelligent agents. Equipped with a rich ecosystem of external tools, agents can tap into domain-specific expertise, access continuously updated real-time data, and interact with users in more flexible ways~\cite{1,3,5}, thereby overcoming the limitations of their inherent knowledge and reasoning capabilities to handle tasks more effectively. 
% 智能体整体性能的瓶颈正在从模型本身转移到其依赖的执行基础设施，也就是 harness。Harness 通过运行时控制提供行为治理与容错能力，使智能体能够在长时任务中保持稳定并在动态环境中自我修正。工具调用是 harness 最核心的治理点之一。已有研究显示，在模型保持不变的情况下，仅调整 harness 中的工具接口格式，就能在 code patch 任务中带来六到十倍的性能差异，由此体现出 harness 对系统表现的决定性影响

% The performance bottleneck of agents is shifting from the large language model (LLM) to the harness layer~\cite{1h,2h,4h}. The harness provides runtime control that enables behavioral governance and fault tolerance, allowing agents to remain stable over complex tasks and adapt to dynamic environments. Research~\cite{3h} shows that tool use is a core focus of the harness: even without changing the agent itself, simply modifying the tool interface can lead to six to ten times better performance in code patching. 
The performance bottleneck of agents is shifting from the capabilities of large language models (LLMs) to the agent harness that controls execution~\cite{1h,2h,4h}.  
Within such harnesses, tools are the primary interface to external environments~\cite{1,5}, and prior studies~\cite{3h} show that minor improvements in tool-use robustness can yield performance gains ranging from 6.7\% to 68.3\%.
% , highlighting the importance of reliable tool use for agent execution.

% Existing tool-use methods still focus on static environments and rely on an implicit assumption that tool-use behavior is predictable. 
% On the one hand, these work still focus on static environments and rely on an implicit assumption that tool-use behavior is predictable. this assumption implies that tools the agent has already mastered can always be invoked correctly and return expected responses. In reality, the quality of tool invocations is strongly affected by environmental dynamics~\citep{7,9}. Even for the same tool, the agent's behavior and answer quality can change with the task context, and may also drift over time or across different usage environments. This issue becomes particularly evident when there is a gap between development and production: many tools that work well in the original setting become unreliable or even unusable because their previous usage patterns no longer apply~\citep{10}.
% On the other hand, the assumption also implies that even when tools are invoked incorrectly, their failures are accompanied by informative signals, enabling the agent to continuously self-reflect and eventually learn to use the tools correctly. In fact, the responses returned by tool invocations are often noisy and coarse-grained~\cite{8}. For instance, a buggy calculator may return an incorrect numerical result while still reporting the invocation as ``successful''. Figure~\ref{fig1}(a) presents examples of these two challenges.

However, the tool interface inside these harnesses rarely adapts across runtime environments. Existing tool-use methods~\cite{toolsurvey} are typically developed under conditions where tool-use behavior remains predictable.
% as they are primarily designed for static situations where tool-use behavior is implicitly assumed to be predictable. 
Concretely, these conditions imply that (1) tools the agent has already mastered can always be invoked correctly and return expected responses, and (2) even when tools are invoked incorrectly, their failures are accompanied by informative signals, enabling the agent to continuously self-reflect and eventually learn to use the tools correctly. 
Unfortunately, real-world conditions break these assumptions. On the one hand, the quality of tool invocations is strongly affected by environmental dynamics~\citep{7, 9}. Even for the same tool, the agent's tool-use behavior can change with the task context. This is particularly evident when there is a gap between development and production: many tools that work well in the original setting become unreliable or even unusable because their previous usage patterns no longer apply~\citep{10}. 
On the other hand, the responses from tool invocations are often noisy or coarse-grained~\cite{8}. For instance, a buggy calculator may return an incorrect numerical result while still reporting the invocation as ``successful''. Figure~\ref{fig1}(a) illustrates examples of these two challenges.

\begin{figure*}[tbp]
    \centering
    \includegraphics[width=0.92\linewidth]{figures/overview.pdf}
    \caption{(a) Two real-world challenges: the agent uses the same tool but gets different results due to the different context. There are different reasons for the error, but the tools respond with the same coarse-grained signal (both return just ``error''). (b) Experience-driven adaptive guidance constructed from past tool use. (c) Results across different stages of tool usage demonstrate that {\ours} improves agent tool-use performance from multiple aspects.}\label{fig1} 
    \vspace{-3mm}
\end{figure*}

% To address these challenges, we propose {\ours} (Dynamic Experience-Driven experience pool), a mechanism that dynamically builds capability boundaries and best practices for tools when agents invoke them blindly. 
% Inspired by agent memory~\cite{memsurvey}, {\ours} acquires, distills, and reuses past experience to produce a unique ``memory'' for each tool, which serves as prior guidance to help the agent interact with the tool more effectively.
% (experiences pool that can be dynamically updated)
To address these challenges, we propose {\ours}, which constructs \textbf{Exp}erience-Driven adaptive \textbf{G}uidance for tool use in the agent harness, so that tool use can improve continually under environmental dynamics and noisy feedback.
% across diverse environments. 
% ensure its. unpredictable
% which serves as prior guidance to help the agent interact with the tool more effectively.
% Specifically, in contrast to existing approaches that rely on static assumptions and treat tool invocations as black-box, isolated, and perfectly
In contrast to existing approaches that rely on static assumptions and treat tool invocations as black-box, isolated, and repeatable events, our work views tool invocations as analyzable, interdependent, and learnable experiences. 
% We partition experiences into equivalence classes aligned with invocation patterns so that guidance remains representative yet covers rare failure modes across acquisition, distillation, and reuse. 
% from past tool invocations
% Inspired by agent memory~\cite{memsurvey}, {\ours} maintains an evolving experience pool by continuously acquiring, distilling, and reusing these experiences, continually building and refining the generated guidance to help agents use tools effectively under environmental dynamics and noisy feedback. 
% By continuously acquiring, distilling, and reusing these experiences, {\ours} maintains an evolving experience pool, inspired by agent memory~\cite{memsurvey}, thus adaptively building and refining guidance for tool use. 
By acquiring, distilling, and reusing tool invocation experiences, {\ours} maintains an evolving experience pool inspired by agent memory~\cite{memsurvey}.
% generalizable
% The guidance characterizes tools' capability boundaries and best practices, helping agents using tools from 
% tool invocation 
Experiences for the same tool are partitioned into equivalence classes after each invocation is attributed to its contributing factors.
% These experiences are partitioned into equivalence classes, such that each class corresponds to a distinct tool invocation pattern. 
% A filtering method is then applied to maintain the quality of the experience pool by removing unhelpful experiences and selecting helpful ones, 
% It builds on the partitioning of the acquired experiences into equivalence classes, 
% thereby retaining representative experiences covering as many invocation patterns as possible while preserving the original distribution. 
% that capture each tool's common success and failure patterns. 
% In contrast to existing approaches that rely on static assumptions and treat tool invocations as black-box, isolated, and perfectly repeatable events, our work views tool invocations as analyzable, interdependent, and learnable experiences. {\ours} acquires these experiences in the first phase by attributing the success and failure of tool invocations from agents' execution trajectories. 
% Our designed selecting method based on equivalence classes is then applied to promote diversity while preserving the original experience distribution, retaining representative experiences that capture each tool's common success and failure patterns. 
Unhelpful experiences are filtered out, and the rest are summarized, through qualitative behavior analysis and quantitative performance analysis, into guidance that states each tool's capability boundaries and best practices.
% The guidance characterizes the capability boundaries and best practices of tools, and can be adaptively refined as the experience pool evolves. 
% which can be continuously refined as more experiences are acquired. 
The guidance is reused whenever the agent invokes the same tool, either as dynamic contextual prompts or as stable schema constraints. As later trajectories enter the pool, this guidance is refined, supporting continually improving tool use in complex and changing applications. 
% enabling more robust and reliable tool use and improving task performance in complex and unpredictable real-world environments.
In summary, our contributions are as follows:
\begin{itemize}%[leftmargin=12pt,topsep=2pt,itemsep=2pt]
    \item  
    % To overcome the challenges posed by reliance on assumption of predictable tool usage,
    % To overcome unpredictable tool-use behaviors caused by environmental dynamics and noisy feedback, 
    % we break out of the static view of tool invocation by characterizing them through a multi-faceted attribution method, thereby acquiring learnable experience that encapsulates comprehensive information about the environment, the agent, and the tool itself.
    We characterize tool invocations by equivalence-class partitioning over different aspects, acquiring experiences that capture the interpretable semantic information of the environment, the agent, and the tool itself.
    % \item We identify two key challenges posed by existing methods' reliance on the assumption of perfect tool invocation environments. We identify the static assumptions commonly relied upon in existing methods for enhancing agent tool use and highlight the challenges these assumptions face in real-world deployment environments. To overcome these challenges, we shift the perspective from viewing tool invocations as black-box, static, and perfectly repeatable isolated events. Instead, we characterize tool invocations by considering the context, the agent, and the tool itself, enabling the acquisition of insightful experiences from agent-tool interactions to improve subsequent tool use. 
    
    \item We propose \textbf{\ours}, which maintains an evolving tool-level experience pool and constructs adaptive guidance for tool use in the agent harness. The guidance captures capability boundaries and best practices and is refined as further experience accumulates. 
    % We propose a dynamic mechanism \textbf{\ours} to maintain memory at the tool level. By acquiring, distilling, and reusing past experiences, {\ours} builds capability boundaries and best practices for tools, which serve as guidance and can be continually refined as the agent becomes more efficient in tool invocation. 
    % builds capability boundaries and best practices for tools, which serve as guidance
    % As the agent continuously interacts with tools, the guidance is dynamically refined, facilitating mutual improvement between the agent and the tools.
    % in the presence of environmental noise, coarse-grained feedback, and imperfect execution
    \item 
    % We conduct extensive experiments covering different stages of tool use, validating the effectiveness of {\ours} across various models, showing that small models equipped with memory can outperform larger memoryless ones. In particular, {\ours} performs notably well under challenging settings, such as noisy environments provided with distractor tools or imperfect information. An in-depth case analysis further demonstrates that {\ours} improves tool-use performance across multiple aspects.
    Extensive experiments across different stages of tool use and model scales show that small models equipped with {\ours} can outperform larger ones without it. Notably, {\ours} remains robust under challenging settings such as distractor tools or imperfect information.
    % An in-depth case analysis further demonstrates that {\ours} improves tool-use performance across multiple aspects.
    % faced with practical challenges. 
    % We conduct extensive experiments in tool-use scenarios characterized by environmental dynamics and noise to systematically evaluate the effectiveness of ReMe. Experimental results demonstrate that ReMe enhances overall performance across the entire process of tool use, significantly improving the agent's robustness and task success rate. 
    % Notably, ReMe exhibits particular advantages in distinguishing between functionally similar or semantically close tools and in correcting unreliable or suboptimal tool invocations, markedly improving the agent's tool-use effectiveness in complex, near-real-world environments.
    % 这种闭环的机制使得智能体能够从历史经验中持续学习，不断优化其工具选择、参数配置和使用策略，最终实现工具使用能力的渐进式提升。
\end{itemize}


\section{Related Work}
We contrast {\ours} with prior work from the two angles of adaptable tool use and a fine-grained harness. Advantages over common method categories are summarized in Appendix~\ref{app:compare}.

\textbf{Tool-augmented LLMs.} 
Tool-augmented LLMs use external tools to extend domain expertise, knowledge access, and interaction~\cite{1,2,5,6,3,4}. Methods that improve tool use either update model parameters or keep them fixed~\cite{toolsurvey}.
% : tuning-based and tuning-free approaches. 
Tuning-based methods internalize tool-use policies by supervised learning~\cite{11,12} or reinforcement learning~\cite{14,15}. Tuning-free methods instead guide a fixed model at inference time, for example by prompts~\cite{16}, chain-of-thought prompting~\cite{17}, or ReAct~\cite{13}. 
These methods already make tools usable. The remaining difficulty is robust use: they are typically developed where tool-use behavior stays predictable~\cite{10,8,7,9}, so a learned invocation pattern does not evolve when the context changes or tool feedback is noisy. 
{\ours} lets tool use continue to improve by constructing guidance from structured invocation experience that captures capability boundaries and best practices. 
% for more effective and reliable tool use
% generation~\citep{89,93,toolsurvey}
% A common strategy is to construct training data containing tool invocations and then perform supervised fine-tuning, enabling the base model to better handle end-to-end tool use. Other work combines reinforcement learning, preference optimization, or multi-stage training, incorporating tool-execution feedback into the training objective to directly improve task success rates and tool-use quality.
% However, both tuning-based and tuning-free methods usually rely on an implicit assumption that tools and environments are relatively stable and predictable. Under this assumption, once the model learns a good invocation pattern, its tool-use performance is expected to transfer smoothly across tasks and deployment settings. In real-world applications, this assumption often breaks down. Tool interfaces, outputs, and underlying data distributions can drift over time and across environments, and tool outputs may be noisy or even contain silent errors.  
% To address this issue, {\ours} explicitly focuses on the robustness of tool use. Our method keeps the flexibility of tuning-free approaches, but adds explicit modeling and accumulation of experience over the entire tool-use process. By extracting, organizing, and reusing experience from past execution trajectories, {\ours} helps agents gradually form a more stable understanding of each tool's capability boundaries and failure modes, and thus use tools more reliably under environmental dynamics and noisy tool feedback.

% \noindent
% \textbf{Learning from Experience.} LLMs have been shown to improve their task-solving capabilities by learning from past experience~\citep{zhao2024expel}, which exhibits particularly strong potential for the continuous evolution of agents.  
% % The core of learning from experience lies in extracting useful information from past interaction trajectories and reusing it in future decision-making and task-solving. to improve its performance
% Early work mostly adopts self-reflection~\cite{18} or self-correction~\cite{19} methods, focusing on analyzing the causes of failures and correcting the agent's subsequent behavior on similar tasks. 
% Recent studies have introduced the concept of memory~\cite{memsurvey}, which goes beyond merely recording and learning from past tasks in a static manner. Instead, memory-based approaches summarize structured knowledge from historical trajectories and build update mechanisms, thereby forming experiences that are both reliable and transferable across different tasks. 
% Mem0~\cite{chhikara2025mem0} is a representative work in this line, enabling agents to interact with the external environment and distill and store key points. However, it relies on prompt-based methods and does not define structured experience for tool invocations to support in-depth analysis. ToolMem~\cite{toolmem} fills this gap by explicitly recording and managing key information involved in tool usage. However, it mainly focuses on the tool selection stage in scenarios where the tools are external generative models, making it difficult to apply to the heterogeneous and complex tool ecosystem encountered in practice. 
% Within this line of work, ToolMem~\cite{toolmem} is the first to establishes an interpretable and maintainable memory framework to the fine-grained tool level, explicitly recording and managing key information involved in tool usage. However, ToolMem mainly focuses on the tool selection stage in a setting scenarios where the tools are external generative models, making it difficult to apply to the heterogeneous and complex tool ecosystem encountered in practice.


\noindent
\textbf{Harness Engineering.} 
Harness engineering improves robustness through the infrastructure around a model rather than through its weights~\cite{18,19,1h,2h}. An execution harness manages the execution loop, tool access, context construction, state persistence, and runtime monitoring~\cite{3h,4h}. Within this infrastructure, memory stores~\cite{memsurvey, chhikara2025mem0}, reusable skills~\cite{1s, 2s}, and context engineering~\cite{c1, c2} externalize knowledge at the task or workflow level. {\ours} revises a tool-specific part of the harness from each tool's own invocation experience. This fine-grained design makes the guidance more specialized and flexible so it can be reused across sessions and environments.

\begin{figure*}[t]
    \centering
    \includegraphics[width=0.93\linewidth]{figures/fmpic.pdf}
    \caption{Overview of {\ours}. It consists of three phases: acquiring tool invocation experiences from past trajectories, distilling them into adaptive guidance, and reusing that guidance in the harness.}\label{fig:overview}
    \vspace{-3mm}
\end{figure*}

% , which can be continuously updated and refined as if the agent has ``memory'' for using the tool effectively.
\section{Methodology}\label{method}
\subsection{Overview of \ours}
Our work, \ours, as illustrated in Figure~\ref{fig:overview}, operates through three phases: experience acquisition, distillation, and reuse. In the \textbf{experience acquisition} phase, an evaluator attributes the success or failure of each tool invocation from historical trajectories, acquiring structured and learnable experiences. In the \textbf{experience distillation} phase, experiences are first filtered to ensure they are helpful, then summarized into guidance that captures the tool's capability boundaries and best practices. Finally, the \textbf{experience reuse} phase applies this guidance during inference, either as contextual prompts or as tool-schema constraints, so that later invocations can improve on earlier experience. 
% These phases work together to build memory for the tool, driving sustained improvements in the agent's performance.
These phases work together to maintain an evolving experience pool for each tool. As further invocations are recorded, the pool and its guidance are updated, so later tool use can improve on earlier experience.
% an evaluator analyzes agent generated trajectories (both successful and failed)
% continuously optimizes the experience pool by incorporating new solid experiences and discarding outdated ones, ensuring stable relevance and adaptability to shifting task demands
\subsection{Experience Acquisition}
We begin by formalizing a tool invocation experience as $E=\langle h,m,c\rangle$, where $h$ is the hash-based unique identifier of the experience. The metadata $m$ captures tool invocation details, including tool information (name and schema), usage information (agent context and tool input), and execution information (tool response, success flag, time cost, and token cost). The experience content $c$ is produced by an agent $\mathrm{LLM}_{\mathrm{evaluator}}$ that takes $m$ as input and consists of: 1) $scores=[s_1,\ldots,s_{d}]$, a $d$-dimensional binary vector with $s_i\in\{0,1\}$ that characterizes different aspects of invocation quality, and 2) $explanation$, a short natural language explanation of the $scores$ that explains why the tool invocation succeeded or failed. 
% and 3) $tool\_problem$, which indicates whether the issue is caused by the tool itself rather than the agent's invocation. 
% tool contexts, inputs, or responses

 % from three different perspectives (tool context, input, and response)
% We design 10 questions (the prompt template is shown in Appendix~\ref{prompt}),
To acquire as many valuable experiences as possible for improving agents' tool use capability, {\ours} first collects trajectories in which an agent solves diverse tasks using different tools. For each agent-tool interaction in these trajectories, metadata $m$ is extracted, with a hash value $h$ computed to distinguish different invocations of the same tool under different tool contexts, inputs, or responses. An agent $\mathrm{LLM}_{\mathrm{evaluator}}$ then evaluates each tool invocation to produce the experience content $c$. 
To ensure accurate and reliable evaluation, our LLM-as-a-Judge~\citep{zheng2023judging} method guides $\mathrm{LLM}_{\mathrm{evaluator}}$ to answer a series of simple yes-or-no questions, such as ``Are the input parameters in the correct format for their respective fields?''. 
The evaluation can be detailed into $d$ aspects, and the answers to the $d$ binary questions are encoded as a $d$-dimensional vector $scores$. 

By providing a natural-language explanation for each dimension, denoted as $explanation$, $\mathrm{LLM}_{\mathrm{evaluator}}$ can attribute the reasons for a tool invocation's success or failure. The prompt template (Appendix~\ref{prompt}) we designed for tool invocation covers 10 aspects, and the corresponding explanation provides comprehensive information by grouping these aspects from three perspectives: (1) \textbf{tool-use environment}, which assesses whether the tool is selected appropriately and invoked at the right time given the context, (2) \textbf{agent behavior}, which assesses whether the agent invokes the tool with an input that is correct both syntactically and semantically, and (3) \textbf{tool reliability}, which assesses whether the tool executes stably and returns a response as expected. All experiences acquired are indexed by tool name and stored in a vector database, forming an initial structured experience pool for the tool. 
% Moreover, we highlight the causes of failures and instruct $\mathrm{LLM}_{\mathrm{evaluator}}$ to judge whether the error is caused by the tool, such as low availability or behavior inconsistent with its schema. 
% multiple
% \footnote{\url{https://en.wikipedia.org/wiki/Equivalence_class}}
With the notion of equivalence classes~\cite{class}, {\ours} naturally utilizes the output of the evaluation and partitions the experience pool $\mathcal{E}_\tau=\{E_i\}_{i=1}^{N_\tau}$ of the same tool $\tau$, into $K$ equivalence classes based on their $scores$, as shown in Eq.~\ref{eq1}: 
\begin{equation}
\label{eq1}
\begin{aligned}
\kappa(E) &\triangleq scores(E)\in\{0,1\}^{d},\\
\{\mathcal{C}_k\}_{k=1}^{K} &\leftarrow \{\kappa^{-1}(v): v\in\kappa(\mathcal{E}_\tau)\}.
\end{aligned}
\end{equation}
Each equivalence class corresponds to a distinct pattern of tool invocation. For example, the perfect invocation pattern can be denoted as class $\mathcal{C}_{\mathrm{perfect}} \triangleq \kappa^{-1}\!\left([1,1,\ldots,1]\right)$, which indicates that the tool is invoked in an appropriate context with correct inputs and returns the expected response. 
% Another special class is defined as $\mathcal{C}_{\mathrm{tool\_problem}} \triangleq \kappa^{-1}\!\left([1,1,1,1,1,1,0,0,0,0]\right)$, which captures cases where the tool is invoked in an appropriate context with correct inputs, yet the response deviates substantially from expectations. In such cases, the failure is more likely to be attributable to the tool itself (e.g., low availability or an implementation inconsistent with its schema) rather than to the agent's invocation behavior.

\subsection{Experience Distillation}\label{sec:reuse}

Equipped with rich experiences, the goal of the second phase in {\ours} is to distill various invocation-level experiences into unified tool-level guidance, which can be flexibly generalized to diverse future tasks and environments. 
Concretely, this phase consists of two steps: \textbf{experience filter} and \textbf{experience summary}. The former keeps the experience pool up-to-date and high-quality by removing unhelpful experiences and selecting helpful ones, while the latter analyzes and summarizes the retained experiences, deriving insightful and practical guidance that characterizes both qualitative tool-use behavior and quantitative tool-use performance. 


% insightful  from comparison and synthesis
% organizing the fragmented experiences into unified and reusable guidance. 

In the experience filtering step, {\ours} removes redundant tool invocation experiences for each tool based on their hash values, but only when the input, output, and context are identical. It then sorts the experiences by their invocation time and discards outdated ones to ensure the timeliness of the experience pool. Then, we perform an effectiveness check with a sliding time window to remove bad experiences followed by sustained degradation. For a given timestamp, if the average $scores$ of the next three consecutive timestamps are consistently lower than that of the current timestamp, all invocation experiences associated with the current timestamp that are used for distilling guidance are removed. 
After removing unhelpful experiences, {\ours} implements different selection strategies by categorizing tools into three types according to the number of experiences: 1) \textbf{one-shot tool}: 
to avoid biased guidance summarized from only a single invocation, no experience 
is selected and the subsequent summary step is skipped for the tool. 2) \textbf{few-shot tool}: all retained experiences are used to summarize comprehensive guidance for the tool. 3) \textbf{many-shot tool}: the equivalence-class-based selection method we designed is applied to select at most $Q$ representative experiences, where $Q$ is computed according to the available context length. 
% We outline the procedure in Algorithm~\ref{algo:ecs_short}, 
The procedure is given as Algorithm~\ref{algo:ecs_short} in Appendix~\ref{proof}, where it is also analyzed in depth. 
The method promotes diversity while preserving the original score distribution, thereby ensuring that the selected experiences are representative by covering as many invocation patterns as possible while emphasizing the most frequently observed ones. 
% as they capture the tool's typical success and failure patterns. 
% selects a maximally diverse subset while preserving the original score distribution as closely as possible
% 在经验过滤（Experience Filter）步骤，{\ours}首先根据每个工具收集的所有经验的hash移除冗余的工具调用经验。然后根据经验的创建时间（往往和工具调用时间有强相关性）进行排序，并移除过期经验，这一点保证了工具记忆的时效性。然后，我们基于时间窗口的对每条经验进行效能判断，用来移除添加后反而造成负面影响的经验。具体来说，{\ours}收集每个时间点该工具的所有经验计算平均得分，若某个时间戳后续连续3个时间戳的平均得分均低于该时间戳，则删除该时间戳下所有已总结的调用记录。接下来，我们精心设计了一种经验选择方法，该方法根据工具的有效经验数将其分成了三类：1）one-shot tool，为了避免从单条调用中获得片面的信息，{\ours}将不选择任何一条经验，使其跳过后一步的总结过程 2）few-shot tool，该工具的调用经验较少，{\ours}将利用所有经验得到全面的guidance，and 3) many-shot tool，该工具的调用经验数量多且各种各样，为此，我们设计了一种基于等价类的选择算法，选择出那些能代表工具使用成功和失败范式的经验，实现如\ref{algo}。算法的核心思想是在保证了所选经验代表性的前提下，尽量保持与原始数据质量分布的一致性，从而为后一步提供全面且准确刻画了工具的调用经验。

After experience filtering is completed, {\ours} conducts qualitative behavior analysis alongside quantitative performance analysis over the retained experiences, producing summarized guidance. The qualitative analysis reflects how the tool should be used. Prompted with chain-of-thought reasoning (see Appendix~\ref{prompt}), an agent $\mathrm{LLM}_{\mathrm{summarizer}}$ can summarize the usage behavior for the tool from four aspects: core function, success patterns, common issues, and best practices. The quantitative analysis reflects the tool's performance and is based on statistics over historical experiences, including the invocations' average success rate, average score, average time cost, and average token cost. 

\subsection{Experience Reuse}\label{sec:refinement}
% accumulated experiences 
The final phase of {\ours} is experience reuse, which applies the constructed guidance to new tasks and shifts the agent from blind trial-and-error tool use to informed tool use. Instead of treating all distilled guidance equally, {\ours} categorizes it into dynamic and stable types based on the consistency of invocation patterns, and applies the two types differently. 
%  {\ours}的最后一个阶段是经验重用，旨在将历史积累的工具使用经验有效应用于新的任务场景中，帮助智能体从试错使用演变为指导推理。 经验重用的核心创新在于一个自适应的应用策略。系统并非简单地将所有检索到的工具guidance统一处理，而是根据历史工具调用结果的统计特征，设计了两种不同的使用模式。具体而言，因为{\ours}是一个动态机制允许工具记忆和生成的guidance都随着时间环境而变化，所以默认的情况下guidance将短期的提示词的形式发挥作用，其中会以较宽松的方式引导模型以一种更智能的方式使用工具。另外一种模式则对应了经过长期积累且经过反复验证的稳定经验，这种时候生成的guidance会转化为工具配置级别的约束条件，以更严格和直接的方式指导模型行为。我们通过公式1来识别这类长期guidance，对应着工具在多次调用中展现出一致的问题模式或被多次标记为工具本身存在问题的情况。值的注意的是，即使该guidance被内化为长期的工具schema，在使用过程中也并非永久不变，随着更多新调用经验的加入，工具记忆可能会被其它使用模式占据主要空间，从而获得新的长期guidance。

Since {\ours} is an adaptive mechanism and the experience pool can be updated over time and across environments, dynamic guidance is by default used as lightweight contextual prompts that guide the agent to invoke tools appropriately. In contrast, stable guidance corresponds to reliable experiences that have been extensively validated. In this case, the distilled guidance is combined into the tool schema, constraining agent behavior more directly and strictly. The categorization criterion is given in Eq.~\ref{eq2}, where $C^{*}$ denotes the largest equivalence class among all classes except the perfect invocation class $\mathcal{C}_{\mathrm{perfect}}$ (perfect tool use does not require stricter constraints), and $\alpha$ is a threshold that controls the strictness of the criterion. Specifically, when the number of experiences reaches half of the selection quota (${Q}/{2}$) and most experiences belong to a single non-perfect class, the criterion function $f$ is triggered and the guidance is classified as stable. We recommend setting $\alpha=K/Q$, following the mixing coefficient of the experience selection strategies analyzed in Appendix~\ref{proof}, which captures the trade-off between diversity coverage and distribution preservation. Using this threshold ensures that stable guidance is triggered only when a high-frequency and consistent equivalence class dominates. The stable guidance covers situations in which a tool exhibits a consistent failure pattern across different invocations. For example, typical stable guidance may reflect tool usability issues in the prompt template. That is, even when the tool is invoked in an appropriate context with correct inputs, its response can still deviate substantially from expectations, causing the equivalence class $\mathcal{C} \triangleq \kappa^{-1}\!\left([1,1,1,1,1,1,0,0,0,0]\right)$ to dominate. 
\begin{equation}
\label{eq2}
% \small
\begin{gathered}
\textsc{Type}(guidance_\tau)
=
\begin{cases}
\text{stable}, & \text{if } f(\mathcal{E}_\tau)\\
\text{dynamic}, & \text{otherwise}
\end{cases}\\[4pt]
f(\mathcal{E}_\tau)
\triangleq
\Big(|\mathcal{E}_\tau|\ge \frac{Q}{2}\Big)\ \wedge\ (|\mathcal{C}^*|\ge \alpha\,|\mathcal{E}_\tau|).
\end{gathered}
\end{equation}

% Another special class is defined as $\mathcal{C}_{\mathrm{tool\_problem}} \triangleq \kappa^{-1}\!\left([1,1,1,1,1,1,0,0,0,0]\right)$, which captures cases where the tool is invoked in an appropriate context with correct inputs, yet the response deviates substantially from expectations. In such cases, the failure is more likely to be attributable to the tool itself (e.g., low availability or an implementation inconsistent with its schema) rather than to the agent's invocation behavior.
% Notably, even after being internalized into a stable tool schema, the guidance is not permanently fixed. As more invocation experiences are acquired, new consistent patterns may become dominant in the experience pool and lead to updated stable guidance.
% \begin{equation}
% \label{eq2}
% \small
% \begin{aligned}
% \textsc{Type}(guidance_\tau)
% &=
% \begin{cases}
% \textsc{stable}, & \text{if } \textsc{LongTerm}(\mathcal{E}_\tau)\\
% \textsc{dynamic}, & \text{otherwise}
% \end{cases}\\[4pt]
% \textsc{LongTerm}(\mathcal{E}_\tau)
% &\triangleq
% \Big(|\mathcal{E}_\tau|\ge \frac{Q}{2}\Big)\ \wedge\ \Big(C_1(\mathcal{E}_\tau)\ \vee\ C_2(\mathcal{E}_\tau)\Big),\\
% C_1(\mathcal{E}_\tau)
% &\triangleq
% |\mathcal{C}^*|\ge \alpha\,|\mathcal{E}_\tau|,\\
% C_2(\mathcal{E}_\tau)
% &\triangleq
% \left|\left\{E\in \mathcal{E}_\tau:\ f(E)\right\}\right|
% \ge \alpha\,|\mathcal{E}_\tau| .
% \end{aligned}
% \end{equation}
% Here, $C^{*}$ denotes the largest equivalence class among all classes except the perfect invocation class
% $\mathcal{C}_{\mathrm{perfect}}$, and $\alpha$ is a threshold that controls the strictness of the criterion. 
% The function $f(E)\in\{0,1\}$ equals 1 if $E\in\mathcal{C}_{\mathrm{tool\_problem}}$ and the
% $LLM_{evaluator}$ attributes the failure of the invocation to \texttt{tool\_problem}.  
Overall, experience reuse leverages distilled guidance to steer tool invocation in new scenarios. These usage strategies allow {\ours} to choose an appropriate way to apply guidance based on its reliability and generalizability. It keeps dynamic guidance flexible to adapt to changing environments, while ensuring the validated stable guidance provides consistent constraints across tasks. As new invocations are recorded, either form of guidance can be revised, so later tool use can improve on earlier experience.


% 这种自适应策略使得系统能够根据经验的成熟度和可靠性，选择最合适的知识传递方式和起作用时间，既能够保持新兴或变化中经验的灵活性以适应各种环境的动态，又能通过配置级约束确保经过充分验证的关键经验得到严格执行，更持续稳定地提升工具使用的成功率和任务整体效能。

\begin{table*}[t]
    \caption{Performance (\%) across MetaTool, API-Bank, and BFCL-V3. \textbf{Bold} indicates the best within each model.}
    \centering
    \fontsize{8.2}{9.5}\selectfont
    \setlength{\tabcolsep}{1.7pt}
    \begin{tabular}{@{}ll*{12}{c}@{}}
    \toprule
         \multirow{2}{*}{\textbf{Model}} & \multirow{2}{*}{\textbf{Method}} &
         \multicolumn{3}{c}{\textbf{MetaTool}} &
         \multicolumn{3}{c}{\textbf{API-Bank}} &
         \multicolumn{3}{c}{\textbf{BFCL-V3}} &
         \multicolumn{3}{c}{\textbf{Total}} \\
         &  & Pass@1 & Avg@3 & Pass@3 & Pass@1 & Avg@3 & Pass@3 & Pass@1 & Avg@3 & Pass@3 & Pass@1 & Avg@3 & Pass@3 \\
    \midrule
    \multirow{5}{*}{\textbf{GPT-5 nano}}
        & No Method & 72.62 & 72.76 & 78.49 & 82.96 & 83.46 & 86.97 & 53.80 & 53.00 & 60.95 & 70.82 & 70.62 & 76.63 \\
        & Few-shot  & 74.12 & 75.11 & 82.32 & 83.71 & 83.96 & \textbf{87.22} & 56.18 & 55.24 & 61.39 & 72.36 & 72.65 & 79.28 \\
        & DRAFT     & 73.94 & 73.04 & 78.97 & 84.21 & 83.46 & \textbf{87.22} & 57.27 & 57.27 & 62.26 & 72.52 & 71.58 & 77.23 \\
        & Mem0      & 74.96 & 76.13 & 82.92 & 84.96 & 85.21 & \textbf{87.22} & 60.95 & 61.61 & 65.08 & 73.98 & 74.67 & 80.35 \\
        & \cellcolor{gray!15}{\ours} & \cellcolor{gray!15}\textbf{81.67} & \cellcolor{gray!15}\textbf{82.07} & \cellcolor{gray!15}\textbf{84.60} & \cellcolor{gray!15}\textbf{86.72} & \cellcolor{gray!15}\textbf{86.55} & \cellcolor{gray!15}\textbf{87.22} & \cellcolor{gray!15}\textbf{64.43} & \cellcolor{gray!15}\textbf{63.99} & \cellcolor{gray!15}\textbf{66.38} & \cellcolor{gray!15}\textbf{79.32} & \cellcolor{gray!15}\textbf{79.22} & \cellcolor{gray!15}\textbf{81.69} \\
    \midrule
    \multirow{5}{*}{\textbf{DeepSeek-V3}}
        & No Method & 83.10 & 82.94 & 84.66 & 84.71 & 84.38 & 85.46 & 58.79 & 59.65 & 65.94 & 78.92 & 78.66 & 81.37 \\
        & Few-shot  & 82.74 & 83.90 & 86.28 & 85.21 & 84.63 & 86.22 & 60.52 & 60.30 & 67.90 & 79.08 & 79.45 & 82.92 \\
        & DRAFT     & 80.23 & 80.79 & 82.44 & 84.96 & 85.63 & 86.47 & 62.26 & 61.61 & 68.55 & 77.70 & 77.80 & 80.54 \\
        & Mem0      & 83.88 & 84.56 & 86.40 & 85.46 & 85.55 & 86.47 & 65.08 & 65.15 & 68.33 & 80.70 & 80.91 & 83.12 \\
        & \cellcolor{gray!15}{\ours} & \cellcolor{gray!15}\textbf{85.26} & \cellcolor{gray!15}\textbf{85.38} & \cellcolor{gray!15}\textbf{86.52} & \cellcolor{gray!15}\textbf{87.72} & \cellcolor{gray!15}\textbf{87.39} & \cellcolor{gray!15}\textbf{87.97} & \cellcolor{gray!15}\textbf{69.41} & \cellcolor{gray!15}\textbf{69.92} & \cellcolor{gray!15}\textbf{72.02} & \cellcolor{gray!15}\textbf{82.76} & \cellcolor{gray!15}\textbf{82.61} & \cellcolor{gray!15}\textbf{84.11} \\
    \midrule
    \multirow{5}{*}{\textbf{Qwen3-8B}}
        & No Method & 76.51 & 76.97 & 77.71 & 83.96 & 83.88 & 84.21 & 58.79 & 58.28 & 60.30 & 74.46 & 74.41 & 75.56 \\
        & Few-shot  & 79.93 & 79.83 & 82.92 & 83.71 & 82.62 & 84.96 & 60.09 & 59.29 & 61.39 & 76.91 & 76.27 & 79.32 \\
        & DRAFT     & 78.19 & 77.33 & 77.89 & 85.71 & 84.96 & 85.46 & 60.74 & 60.30 & 62.91 & 76.20 & 75.18 & 76.35 \\
        & Mem0      & 75.07 & 75.47 & 82.38 & 86.22 & 86.05 & 86.47 & 63.34 & 64.93 & 66.16 & 74.69 & 74.98 & 80.07 \\
        & \cellcolor{gray!15}{\ours} & \cellcolor{gray!15}\textbf{83.52} & \cellcolor{gray!15}\textbf{84.88} & \cellcolor{gray!15}\textbf{85.08} & \cellcolor{gray!15}\textbf{86.47} & \cellcolor{gray!15}\textbf{87.89} & \cellcolor{gray!15}\textbf{87.97} & \cellcolor{gray!15}\textbf{67.46} & \cellcolor{gray!15}\textbf{66.96} & \cellcolor{gray!15}\textbf{68.33} & \cellcolor{gray!15}\textbf{81.06} & \cellcolor{gray!15}\textbf{81.82} & \cellcolor{gray!15}\textbf{82.48} \\
    \midrule
    \multirow{2}{*}{\textbf{Qwen3-32B}}
        & No Method & 80.05 & 79.43 & 80.17 & 84.71 & 84.88 & 85.21 & 65.15 & 65.08 & 66.16 & 78.05 & 77.55 & 78.41 \\
        & \cellcolor{gray!15}{\ours} & \cellcolor{gray!15}\textbf{84.68} & \cellcolor{gray!15}\textbf{85.02} & \cellcolor{gray!15}\textbf{86.28} & \cellcolor{gray!15}\textbf{86.97} & \cellcolor{gray!15}\textbf{87.30} & \cellcolor{gray!15}\textbf{87.72} & \cellcolor{gray!15}\textbf{70.72} & \cellcolor{gray!15}\textbf{71.01} & \cellcolor{gray!15}\textbf{73.32} & \cellcolor{gray!15}\textbf{82.48} & \cellcolor{gray!15}\textbf{82.56} & \cellcolor{gray!15}\textbf{84.14} \\
    \midrule
    \multirow{2}{*}{\textbf{Qwen3-235B}}
        & No Method & 78.25 & 79.23 & 80.29 & 85.46 & 85.46 & 85.71 & 71.37 & 71.15 & 73.54 & 78.13 & 78.49 & 79.91 \\
        & \cellcolor{gray!15}{\ours} & \cellcolor{gray!15}\textbf{86.34} & \cellcolor{gray!15}\textbf{86.70} & \cellcolor{gray!15}\textbf{86.94} & \cellcolor{gray!15}\textbf{87.47} & \cellcolor{gray!15}\textbf{86.97} & \cellcolor{gray!15}\textbf{88.22} & \cellcolor{gray!15}\textbf{79.61} & \cellcolor{gray!15}\textbf{78.52} & \cellcolor{gray!15}\textbf{80.04} & \cellcolor{gray!15}\textbf{85.29} & \cellcolor{gray!15}\textbf{84.98} & \cellcolor{gray!15}\textbf{85.69} \\
    \bottomrule
    \end{tabular}
    \label{tab:main}
% \vspace{-1mm}
\end{table*}

\section{Experiments}
\subsection{Experimental Settings}

\paragraph{Baselines.} 
To evaluate the effectiveness of {\ours}, we compare it with: (1) No Method, and three types of methods that improve the agent's tool use by learning from prior experience: (2) Few-shot learning~\cite{fewshot}, a prompting method that provides a small number of tool-use examples in the prompt to let the model learn from them, (3) DRAFT~\cite{Draft}, a method that dynamically improves the quality of tool documents to enhance agents' understanding and use of tools, and (4) Mem0~\citep{chhikara2025mem0}, a method which builds conversational memory for an agent by extracting key information from its interactions with the external world. 
These three methods provide guidance for tool-use from different aspects, with further details provided in Appendix~\ref{app:baselines}. 
% (2)Dspy~\cite{dspy}, a method for improving an agent's prompts by creating and collecting examples from previously completed tasks

\paragraph{Datasets and Metrics.} % Datasets and Metrics
The task-solving process of tool-augmented LLMs is often decomposed into four stages: task planning, tool selection, tool calling, and response generation~\cite{toolsurvey}.
To thoroughly evaluate the effectiveness of {\ours}, we conduct experiments on three benchmark datasets, each targeting different stages of tool usage: MetaTool~\cite{metatool}, API-Bank~\cite{apibank}, BFCL-V3~\citep{patilberkeley}. 
For MetaTool, we select the 100 most frequently used tools and sample up to 10 tasks for each tool under each sub-task setting, resulting in 1,669 tasks in total. For API-Bank, we use the full Level-1 dataset, comprising 399 tasks. Both datasets are evaluated using leave-one-out cross-validation. For BFCL-V3, since some tasks involve tools that are never reused in other tasks, we filter out these cases and retain 461 tasks for the test set and 1,542 tasks for the training set. 
We use both Avg@3 and Pass@3 as evaluation metrics. Avg@3 denotes the average task success rate over three independent runs, while Pass@3 measures the probability that at least one of the three runs is successful. As different datasets target different stages of tool usage, we follow their original evaluation criteria when defining ``successful'' to maintain consistency. See Appendix~\ref{app:dataset} for more dataset and metric details.

\paragraph{Implementation Details.} 
We conduct experiments on GPT-5 nano~\cite{openai2025gpt5}, DeepSeek-V3~\cite{deepseekai2025deepseekv3technicalreport}, and the Qwen3 series models~\citep{qwen3technicalreport} to demonstrate the effectiveness of {\ours} across different model architectures and scales. We use a prompt-based tool invocation method and set $temperature =0.5$ for trajectory sampling. 
During experience acquisition and experience distillation, we adopt Qwen3-max as the $\mathrm{LLM}_{\mathrm{evaluator}}$ and $\mathrm{LLM}_{\mathrm{summarizer}}$ to produce reliable tool invocation experiences and distilled guidance. 
During experience reuse, we set $\alpha=K/Q$ to control the criterion in this usage strategy for guidance. 
In addition, to support our design of removing unhelpful experiences over time, we execute tasks sequentially. 
For fair comparison, we keep these settings fixed unless otherwise specified in ablation studies. 
% including the extraction process of metadata $m$ in the experiences， the automatic calculation method for the hyperparameter quota $Q$, and other things we tried
More implementation details can be found in the Appendix~\ref{app:tried}. 


\subsection{Main Results} 
% 表~\ref{tab:results} 汇总了五种骨干模型在 MetaTool、API-Bank 与 BFCL-V3 上的实验结果。在每个模型组中，{\ours} 都在三项基准的 Avg@3 与 Pass@3 上取得最优成绩，从而带来sota的表现。以 Qwen3-8B 为例，相比 No Method，{\ours} 在Total 指标上分别带来+7.41 的 Avg@3 与 +6.92 的 Pass@3，且这种增益可以持续到更大规模的模型变体上，验证了我们方法的通用性。值得注意的是，配备 {\ours} 的小模型能够达到与更大模型（但不使用记忆）相当甚至更优的表现，这表明，有效的记忆机制能够显著缩小不同模型规模之间的性能差距。
Table~\ref{tab:main} reports the main results of {\ours} across models on three datasets that target different stages of tool invocation. 
Overall, {\ours} achieves the best Avg@3 and Pass@3 on every dataset, consistently outperforming all competing methods.   
For Qwen3-8B, {\ours} improves the total metric by
{\setlength{\fboxsep}{0.5pt}\colorbox{gainbg}{7.41$\uparrow$}} in Avg@3 and
{\setlength{\fboxsep}{0.5pt}\colorbox{gainbg}{6.92$\uparrow$}} in Pass@3 over the ``No Method'' setting. We also report Pass@1, which highlights clearer distinctions between methods under lower inference budgets. 
These gains persist as model size increases, demonstrating the generality of our approach. 
Importantly, {\ours} allows smaller models to match or even surpass larger models in the ``No Method'' setting, indicating that our approach helps mitigate the tool-use capability gap between weaker and stronger models. 

\begin{figure}[t]
    \centering
    \includegraphics[width=\linewidth]{figures/output.png}
    \caption{Avg@3 performance (\%) on MetaTool and BFCL-V3 across sub-task settings.}
    \label{fig:detail}
    \vspace{-5mm}
\end{figure}

% 除总体表现外，我们还对每个基准进行了深入分析，图1给出了在MetaTool和BFCL-V3数据集中的子任务上的表现。具体来说，工具选择（MetaTool）任务中，{\ours} 在区分相似工具和工具存在可靠性问题时格外突出。工具调用（API-Bank）任务中，我们观察到了不同模型在几个方法上的相似指标，并推测是由于该数据集设定相对简单，到达了针对工具调用方面的优化上限所造成的。在BFCL-V3上，我们的方法在问题设定（缺失参数，缺失方法，不相关方法等）子集上的表现非常亮眼，这表明{\ours}从过去经验中学习到了如何正确面对这些不完美情况。此外，一个有趣的统计是我们的方法在常规的多轮任务中平均推理步骤从10.14减少至9.02，而针对有问题的多轮任务则增加了2.4步。这表明正常推理时{\ours}帮助模型在较少次数的尝试内提高成功概率，而面对不完美情况时候，模型有意识通过和用户交互或提前检查环境去解决这些问题。
Beyond the overall performance, we further conduct a detailed analysis of {\ours} on each dataset. 
Figure~\ref{fig:detail} presents the results on MetaTool and BFCL-V3 across different sub-task settings. 
In tool selection (MetaTool) tasks, {\ours} is particularly effective ({\setlength{\fboxsep}{0.5pt}\colorbox{gainbg}{16.34$\uparrow$}}) when faced with distractor tools that are very similar. 
% or have reliability issues({\setlength{\fboxsep}{0.5pt}\colorbox{gainbg}{6.20$\uparrow$}}). the model is
In tool calling (API-Bank) tasks, we observe that different methods achieve comparable scores. 
% tool-augmented
This is likely because the dataset is relatively simple, allowing all approaches to bring model performance close to the upper limit of their inherent capabilities. 
In response generation tasks (BFCL-V3), {\ours} demonstrates strong performance ({\setlength{\fboxsep}{0.5pt}\colorbox{gainbg}{13.09$\uparrow$}} in single-turn tasks and {\setlength{\fboxsep}{0.5pt}\colorbox{gainbg}{10.33$\uparrow$}} in multi-turn tasks) on noisy settings, such as missing parameters, undefined tools or irrelevant tool choices. 
% on imperfect information settings
% This indicates that {\ours} helps the agent generalize tool use effectively in unpredictable situations. 
This indicates that {\ours} helps the agent generalize tool use effectively in unpredictable or even problematic environments.
% This indicates {\ours} help agent successfully generalizes the tool usage to imperfect scenarios. problematic
% This suggests that the agent has learned to handle such imperfect scenarios based on prior experience. 
Besides, we observe an interesting phenomenon that our method reduces the average inference steps in normal multi-turn tasks from 10.14 to 9.02, but increases the step count by 2.4 in problematic multi-turn tasks. 
% This indicates that {\ours} enables the model to achieve higher success rates with fewer trials under normal settings, while encouraging more deliberate behaviors when facing challenging scenarios. 
This indicates that {\ours} enables the agent to achieve higher success rates with fewer trials under normal settings, reaching correct solutions with fewer unnecessary interactions. At the same time, it encourages more deliberate behaviors in challenging scenarios, prompting the agent to interact with the environment to ensure more reliable tool use. 
Such a shift toward a more flexible tool-use strategy is desirable, as it suggests that the agent becomes more robust across diverse situations under our guidance.
% by consciously interacting with users or conducting anticipatory environmental checks to address potential problems.



\begin{figure}[t]
    \centering
    \includegraphics[width=0.95\linewidth]{figures/ablation.png}
    \caption{Ablation of different phases (Avg@3, \%).}
    \label{fig:ablation}
    \vspace{-5mm}
\end{figure}

\subsection{Ablation Studies}\label{sec:ablation}
As shown in Figure~\ref{fig:ablation}, we conduct an ablation study of the key components of {\ours} using Qwen3-8B as the backbone model. Compared with the baseline without any method, progressively adding each phase yields consistent improvements. First, only using \textit{experience acquisition} as guidance during inference achieves better performance than few-shot learning, indicating that structured experiences acquired from task interactions provide more effective signals. Second, adding \textit{experience distillation} results in a substantial additional gain, highlighting the importance of maintaining a high-quality experience pool and aggregating multiple experiences. Finally, although the additional improvement after adding \textit{experience reuse} is milder than that of the first two phases, it is crucial for driving long-term evolution of the agent and tool quality in real applications. We also tried multiple variants within phases and report the results in Appendix~\ref{app:variants}. 


\begin{table}[tbp]
\caption{Effect of the number of acquired experiences on MetaTool performance (Avg@3, \%) in {\ours}. $\Delta$ denotes the absolute change compared to the previous setting.}
\vspace{-1mm}
\centering
\fontsize{8}{9}\selectfont
\setlength{\tabcolsep}{1.5pt}
\begin{tabular*}{\columnwidth}{@{\extracolsep{\fill}}l*{8}{c}@{}}
\toprule
\textbf{Count} & 0 & 3 & 6 & 9 & 11 & 14 & 17 & 20 \\
\midrule
\textbf{Avg@3} & 77.29 & 82.26 & 85.80 & 86.70 & 86.16 & 84.42 & 85.20 & 86.22 \\
\textbf{$\Delta$} & -- & +4.97 & +3.54 & +0.90 & $-0.54$ & $-1.74$ & +0.78 & +1.02 \\
\bottomrule
\end{tabular*}
\label{tab:experience_num}
\vspace{-3mm}
\end{table}




Additionally, to study how the number of experiences affects {\ours}, we conduct experiments on MetaTool, selecting tools with sufficiently diverse invocations. Table~\ref{tab:experience_num} shows that {\ours} improves quickly once the number of experiences for a tool meets our few-shot tool criterion. As more experiences are acquired, performance peaks at around 9 experiences, after which adding additional experiences introduces fluctuations. With further increases in experience, {\ours} stabilizes and maintains strong performance, as the experience filtering phase continuously removes unhelpful experiences while keeping the most helpful ones. 



\subsection{More Analysis}\label{sec:moreana}

\paragraph{Agent Capabilities.} Since {\ours} relies on the LLM-as-a-Judge method to acquire experiences and summarize guidance effectively, we examine the impact of agent size for both roles, ${\text{LLM}}_{\text{evaluator}}$ and ${\text{LLM}}_{\text{summarizer}}$. The experiments are conducted on BFCL-V3, and the results are shown in Table~\ref{tab:summarizer}. Replacing either agent with a smaller model degrades the final task performance, with a more pronounced drop observed for ${\text{LLM}}_{\text{evaluator}}$. Specifically, substituting the ${\text{LLM}}_{\text{evaluator}}$ with the weaker Qwen3-8B yields a larger performance decline, likely because processing complex trajectories and making accurate attributions is challenging for weaker agents. In contrast, using a moderately sized model such as Qwen3-32B leads to a relatively smaller degradation. 
To further check whether the gains mainly come from a high-tier model, we replace both the evaluator and the summarizer with Qwen3-8B. Table~\ref{tab:same_scale} shows that even without a high-tier model, {\ours} still yields large and stable gains over No Method on all three datasets, and these gains remain comparable to those of Qwen3-max. This indicates that the gains mainly come from structured distillation and reuse of tool experiences, and that agents can self-improve through {\ours} without a stronger external model.

\begin{table*}[t]
\centering
\begin{minipage}{0.48\linewidth}
\centering
\caption{Impact of agent size on BFCL-V3 performance (\%) in {\ours}. $\triangle$ indicates the absolute change relative to the original Qwen3-max setting.}
\vspace{-0.5mm}
\footnotesize
\setlength\tabcolsep{2pt}
\scalebox{0.84}{%
    \begin{tabular}{c|c|cc}
    \toprule
    \multirow{2}{*}{\textbf{${\text{LLM}}_{\text{evaluator}}$}} &
    \multirow{2}{*}{\textbf{${\text{LLM}}_{\text{summarizer}}$}} &
    \multicolumn{2}{c}{\textbf{BFCL-V3}}\\
    && Avg@3 (\%)& Pass@3 (\%)\\
    \midrule
    \multirow{2}{*}{Qwen3-max} & Qwen3-8B & 63.20\colorbox{pink!20}{$\textcolor{darkred}{\triangle=3.76 \downarrow}$} & 65.72\colorbox{pink!20} {$\textcolor{darkred}{\triangle=2.61 \downarrow}$} \\
    &Qwen3-32B & 65.51 \colorbox{pink!20}{$\textcolor{darkred}{\triangle=1.45 \downarrow}$} & 66.02 \colorbox{pink!20}{$\textcolor{darkred}{\triangle=2.31 \downarrow}$} \\
    \midrule
    Qwen3-8B & \multirow{2}{*}{Qwen3-max} &
    62.18\colorbox{pink!20}{$\textcolor{darkred}{\triangle=4.78 \downarrow}$} & 64.71\colorbox{pink!20} {$\textcolor{darkred}{\triangle=3.62 \downarrow}$} \\
    Qwen3-32B &  &
    64.78\colorbox{pink!20}{$\textcolor{darkred}{\triangle=2.18 \downarrow}$} &
    66.81\colorbox{pink!20}{$\textcolor{darkred}{\triangle=1.52 \downarrow}$} \\
    \midrule
    Qwen3-max & Qwen3-max &66.96 & 68.33 \\
    \bottomrule
    \end{tabular}%
}
\label{tab:summarizer}
\end{minipage}\hfill
\begin{minipage}{0.48\linewidth}
\centering
\caption{Self-improving {\ours} with Qwen3-8B (\%). $\triangle$ indicates the absolute change relative to No Method.}
\vspace{-1mm}
\footnotesize
\setlength\tabcolsep{1.2pt}
\scalebox{0.83}{%
    \begin{tabular}{cc@{\hspace{6pt}}|@{\hspace{4pt}}ccc}
    \toprule
    \multirow{2}{*}{\textbf{Dataset}} &
    \multirow{2}{*}{\textbf{Metric}} &
    \multirow{2}{*}{\textbf{No Method}} &
    \multicolumn{2}{c}{\textbf{\ours}} \\
    &&& Qwen3-8B & Qwen3-max \\
    \midrule
    \multirow{2}{*}{MetaTool} & Avg@3 & 76.97 & 81.91\colorbox{gainbg}{$\textcolor{darkgreen}{\triangle=4.94 \uparrow}$} & 84.88\colorbox{gainbg}{$\textcolor{darkgreen}{\triangle=7.91 \uparrow}$} \\
    & Pass@3 & 77.71 & 82.03\colorbox{gainbg}{$\textcolor{darkgreen}{\triangle=4.32 \uparrow}$} & 85.08\colorbox{gainbg}{$\textcolor{darkgreen}{\triangle=7.37 \uparrow}$} \\
    \midrule
    \multirow{2}{*}{API-Bank} & Avg@3 & 83.88 & 85.88\colorbox{gainbg}{$\textcolor{darkgreen}{\triangle=2.00 \uparrow}$} & 87.89\colorbox{gainbg}{$\textcolor{darkgreen}{\triangle=4.01 \uparrow}$} \\
    & Pass@3 & 84.21 & 85.96\colorbox{gainbg}{$\textcolor{darkgreen}{\triangle=1.75 \uparrow}$} & 87.97\colorbox{gainbg}{$\textcolor{darkgreen}{\triangle=3.76 \uparrow}$} \\
    \midrule
    \multirow{2}{*}{BFCL-V3} & Avg@3 & 58.28 & 62.18\colorbox{gainbg}{$\textcolor{darkgreen}{\triangle=3.90 \uparrow}$} & 66.96\colorbox{gainbg}{$\textcolor{darkgreen}{\triangle=8.68 \uparrow}$} \\
    & Pass@3 & 60.30 & 63.99\colorbox{gainbg}{$\textcolor{darkgreen}{\triangle=3.69 \uparrow}$} & 68.33\colorbox{gainbg}{$\textcolor{darkgreen}{\triangle=8.03 \uparrow}$} \\
    \bottomrule
    \end{tabular}%
}
\label{tab:same_scale}
\end{minipage}
\end{table*}

\begin{figure*}[tb]
    \centering
    \includegraphics[width=0.8\linewidth]{figures/sampling_comparison_3d_d_q.png}
    \caption{Distribution alignment and diversity under different selection methods.}
    \label{fig:selection}
    \vspace{-3mm}
\end{figure*}

\paragraph{Equivalence-Class-Based Selection Method.} We compare our equivalence-class-based experience selection method with two typical selection methods: instance-balanced, which selects experiences uniformly, and class-balanced, which selects similar numbers from each class. We use negative Kullback-Leibler divergence~\cite{kl} to measure distribution alignment with the original experience distribution, and Shannon entropy~\cite{entropy} to measure the diversity of the selected set. Figure~\ref{fig:selection} shows the results under different dimensions $d$ of the $scores$ and selection ratios. In terms of distribution alignment, our method performs similarly to the instance-balanced method, maintaining strong consistency with the original distribution. It can even outperform the instance-balanced method through fine-grained adjustment of the selection quotas at higher selection ratios. From the perspective of diversity, the class-balanced method usually achieves the highest entropy by equalizing class proportions, whereas the instance-balanced method has lower entropy, indicating weaker class coverage. Our method achieves much higher entropy than the instance-balanced method and slightly lower entropy than the class-balanced method in most settings. It ensures that the chosen experiences both cover as many invocation patterns as possible and remain representative of the original distribution.


% 我们将经验选择步骤中基于等价类的选择方法其与两个经典选择方法进行对比。Instance-balanced方法对每个样本一视同仁、按实例均匀无放回采样。Class-balanced方法让每个类别被选中的次数尽量相近。图~\ref{fig:selection} illustrates how the three methods behave across different dimensions $d$ in $scores$ and selection ratios. 我们使用负kl来衡量其与原始经验分布的distribution alignment，使用entropy来衡量选择后经验集合的diversity。在一致性（distribution alignment）角度来看，我们的方法有着与 Instance-balanced 相近的表现，保持着与原始分布较高的一致性。此外，在高选择比例中，因其第二阶段细粒度的调整，甚至得到了比Instance-balanced更优的表现。从多样性（diversity）角度来看，Class-balanced 在通常获得最高的 entropy，因为它强调最大化类别覆盖。 Instance-balanced则体现出较低的值，说明它对稀有模式的经验覆盖不足。我们的Equivalence-class-based 方法在几乎所有设置下显著高于 Instance-balanced 而略低于 Class-balanced方法。它通过在等价类层面进行控制，从而确保所选择的经验既覆盖了尽可能多的调用模式同时又具有原始分布代表性。

% 有效提升了所选经验的模式多样性，体现出在保持分布代表性的前提下显著增强多样性的优势。
% 在强调等价类覆盖的同时，并没有显著破坏原始分布的一致性，
% Instance-balanced始终保持着接近原始经验分布的分布while Class-balanced因其拉平各个类别的概率表现始终较差。
% ，说明其在强调等价类覆盖的同时，并没有显著破坏原始分布的一致性
% Equivalence-class-based（我们的方法）：先在等价类层面建模，再在类内按实例采样，在保持整体分布对齐的同时，有针对性地覆盖更多代表性等价类，在分布对齐和多样性之间取得更好的平衡。
% 本文概述了算法\ref{algo:ecs_short}中的过程，在保留原始分数分布的同时促进多样性，通过覆盖尽可能多的调用模式同时强调最频繁观察到的模式，从而确保所选择的经验具有代表性。

\vspace{-1mm}
\paragraph{Case Analysis.} We perform case analysis to better understand the strengths and limitations of {\ours}. These cases show that {\ours} mainly helps the agent: 1) improve tool-use awareness, 2) distinguish similar tools, 3) invoke tools correctly, 4) fix tool-use errors, 5) identify tool dependencies, and 6) understand tool reliability. We also study cases where {\ours} fails, such as unseen situations, hard-to-analyze trajectories, incorrect attributions, and over-reaction. Detailed analyses and a case-coverage comparison with other methods are in Appendix~\ref{case}.

\section{Conclusion}
% generalizable capability boundaries and best practices
% This paper aims to improve the robustness of tool-augmented LLM agents in unpredictable real-world settings.
\vspace{-1mm}
This paper studies how to construct experience-driven adaptive guidance for continually improving tool use in agent harnesses. We propose {\ours}, which acquires, distills, and reuses tool invocation experience, then applies the resulting guidance as contextual prompts or schema constraints. Experiments on tool selection, tool calling, and response generation show that this guidance makes tool use more reliable, and the same mechanism can refine the guidance as further experience is incorporated. We believe {\ours} is a step toward tool use that improves continually as experience accumulates.
% We believe {\ours} represents an early step toward trainable tools, paving the way for building continuously self-improving agent systems in the future.


\section*{Limitations}
Several limitations remain.
First, generating high-quality guidance depends on the evaluator and summarizer agents. Each tool invocation triggers two extra agent inferences as well as storage and retrieval costs, which incurs additional time and cost.
Second, {\ours} requires initial interaction trajectories. If a tool is never invoked, or if some success or failure patterns are not encountered early on, it may lack the experiences needed to form guidance.

\bibliography{expg}
%%%%%%%%%%%%%%%%%%%%%%%%%%%%%%%%%%%%%%%%%%%%%%%%%%%%%%%%%%%%

\appendix

\section{Detailed Comparison with Existing Work}\label{app:compare}
This appendix states the advantage of {\ours} over common categories of prior methods.

\textbf{Tuning-based methods} learn a tool-use policy in the model weights by supervised learning or reinforcement learning. {\ours} saves that training cost and is more flexible: the weights stay fixed, and new invocation experience updates the tool-level guidance in the harness.

\textbf{Inference-time reasoning methods} guide the current task with a fixed procedure or a static tool description. They are typically developed where tool-use behavior stays predictable. {\ours} accumulates invocation experience and retains it for later tasks. In these methods, the experience of one task is not kept as lasting guidance.

\textbf{RAG and few-shot learning methods} place external context in the prompt. RAG returns stored text chunks to cover missing or outdated facts. Few-shot learning, as a prompt cache, concatenates raw tool-call examples. {\ours} gives later tasks refined guidance on capability boundaries and best practices. Unhelpful experiences are removed, and new invocations update the guidance those tasks read.

\textbf{Harness-related methods} externalize knowledge for a task or a workflow, through memory, reusable skills, or context construction. {\ours} is more specialized: the guidance is fine-grained and written for one tool, so it stays with that tool across tasks, sessions, and environments. 

\section{Template Prompt for LLM-as-a-Judge Agent}\label{prompt}

\begin{tcolorbox}[title={Prompt for agent ${\text{LLM}}_{\text{evaluator}}$},breakable,]
You are an expert in evaluating tool invocation process. The tool is invoked by an AI agent.\\

\textbf{Tool invocation Information:}\\
$\bullet$ Tool Name: \{tool\_name\}\\
$\bullet$ Success Flag: \{success\_flag\}\\
$\bullet$ Time Cost: \{time\_cost\}s\\
$\bullet$ Token Cost: \{token\_cost\} tokens\\
$\bullet$ Agent Context: \{context\}\\
$\bullet$ Input Parameters: \{input\_params\}\\
$\bullet$ Tool Response: \{response\}\\
$\bullet$ Tool Schema: \{schema\}\\

\textbf{Evaluation Method:} Start from a default score list of scores = [0, 0, 0, 0, 0, 0, 0, 0, 0, 0]. For each item below that is satisfied, assign 1 point to the corresponding index. The final scores should be a list of 10 integers, each being either 0 or 1.\\

\textbf{1. Tool-use environment (total 2 points. If context is provided, use it as an aid when evaluating):}\\
$\bullet$ Index 1: Should the tool be invoked now? Consider whether the all necessary information for the tool's invocation is ready, and whether the tool execution environment is correct. If it is a multi-round conversation, also consider the dependency relationships of the tool chain.\\
$\bullet$ Index 2: If should, is the chosen tool appropriate?\\

\textbf{2. Agent behavior (total 4 points. When evaluating, consider both the context and the tool schema):}\\
$\bullet$ Index 3: Are all required parameters provided?\\
$\bullet$ Index 4: Are the input parameters valid and supported by the tool?\\
$\bullet$ Index 5: Are the input parameters in the correct format for their respective fields?\\
$\bullet$ Index 6: Does the value(content) of input parameter correctly reflect and match the given context?\\

\textbf{3. Tool reliability (total 4 points):}\\
$\bullet$ Index 7: Does the response provide meaningful and useful information? or are there any error messages or information can be used as guidance for agent invoking tool better?\\
$\bullet$ Index 8: Does the response match the tool's intended purpose/function?\\
$\bullet$ Index 9: Does the response value correct (content appropriate) given the input parameters?\\
$\bullet$ Index 10: Does the response help accomplish the task within the given context?\\

\textbf{Important:}\\
1. Sometimes there is not enough information in the context or schema to make a complete evaluation. In such cases, make your best judgment based on the available information.\\
2. Some tools (commonly system tools such as mkdir, touch, echo, etc.) modify the external environment. Since these results cannot be obtained, they return ``None'' as the response. At this point, all the scores in the quality of the response should be obtained and should not be seen as a problem for the tool.\\
3. Evaluation independently from the success flag. The \texttt{success\_flag} indicates whether the tool executed without technical errors. The \texttt{evaluation} should evaluate the quality of the tool invocation. A tool can execute successfully (Success Flag=1) but still produce low-quality or irrelevant responses, leading to a low evaluation score.\\
4. Sometimes an agent will execute multiple steps and invoke multiple tools to complete a task, but you only need to evaluate the use of one tool for one of the steps, not whether the final task is completed or not.\\

\textbf{Answer Format:}\\
Please provide your answer in the following JSON format:\\

\verb|```|json\\
\{\\
\quad ``scores'': A list of 10 integers (0 or 1) based on the evaluation criteria above.\\ 
\quad ``explanation'': A brief evaluation (2-3 sentences) explaining the quality of the tool invocation, based on your evaluation. Low-quality aspects need to be reified, especially the causes of tool invocation errors.\\
\}\\
\verb|```|
\end{tcolorbox}

\begin{tcolorbox}[title={Prompt for agent ${\text{LLM}}_{\text{summarizer}}$},breakable,]
You are an expert in analyzing tool usage patterns and generating practical usage guidance for agents.\\

\textbf{Tool Information:}\\
$\bullet$ Tool Name: \{tool\_name\}\\
$\bullet$ Tool Schema: \{tool\_schema\}\\

\textbf{Recent Tool Invocation Experiences:}\\
\{experiences\}\\

\textbf{Important:}\\
1. Assume the tool (tool schema) can't be changed, your task is to guide agent to use it better.\\
2. Your answer must be based on the information given, don't make it up. If not enough data, state ``Not enough data to determine Core Function/Success Patterns/Common Issues/Best Practices.''\\
3. Your answer will be used to guide the use of the tool in the future, so do not include content related to recent tool invocation experience such as ``case \#3'' or ``Call \#2'', but some values can be used as examples.\\
4. Pay attention to information not mentioned in the tool schema, such as the response upon successful tool invocation. It's also welcome to uncover insights, such as how tools can be used more effectively, and possible dependencies between tools. But if they aren't, don't make them up.\\
5. Finally, to avoid deriving incorrect guidance from individual invocation, check whether, if the agent follows the proposed guidance, it can perform better on all recent invocation histories. If not, revise the guidance until it can. Specifically:\\
\quad$\bullet$ don't write guidance in an absolute tone without a very deterministic message (meaning that all invocation histories are satisfied, otherwise it will result in failure).\\
\quad$\bullet$ sometimes there may be inconsistencies. Consider whether this is due to the context in which the tool is being used.\\

\textbf{Your Task:}\\
Based on the tool invocation history, generate a concise and logical tool usage guidance following this structure:\\
1. \textbf{Core Function}: What this tool does and when to use it.\\
2. \textbf{Success Patterns}: Parameter patterns and usage scenarios that work well.\\
3. \textbf{Common Issues}: Main pitfalls to avoid and why they fail.\\
4. \textbf{Best Practices}: 2-3 actionable recommendations.\\

\textbf{Answer Format:}\\
Provide a structured, concise guidance (max 200 words). Focus on actionable insights derived from actual usage data. Avoid generic advice and think step by step.\\

\verb|```|txt\\
Your concise, data-driven tool usage guidance\\
\verb|```|
\end{tcolorbox}


% \textcolor{red}{Please skip this section while reading it and revise it after we decided on the algorithm.}
\section{Analysis of Equivalence-Class-Based Selection Method}\label{proof}

{
\SetAlgoSkip{1pt}
\begin{algorithm}[t]
\SetAlCapFnt{\small}
\SetAlCapNameFnt{\small}
\DontPrintSemicolon
\caption{Equivalence-Class-Based Experience Selection Algorithm for Many-shot Tool}
\label{algo:ecs_short}
\small
\KwIn{Equivalence classes $\{\mathcal{C}_k\}_{k=1}^{K}$ of tool $\tau$ derived from $\mathcal{E}_\tau$, and target selection quota $Q$}
\KwOut{Selected subset $\mathcal{S}_\tau$}

\BlankLine
{\color{lightblue}{// Select one from each of the $Q$ largest classes when the number of classes exceeds the quota}}\;
\If{$K>Q$}{
  $(\mathcal{C}_{(1)},\ldots,\mathcal{C}_{(K)}) \leftarrow \textsc{Sort}_{\mathrm{desc}}(\{\mathcal{C}_k\}_{k=1}^{K}, \,|\cdot|)$\;
  $\mathcal{S}_\tau \leftarrow \emptyset$\;
  \For{$k\leftarrow 1$ \KwTo $Q$}{
    $E \leftarrow \textsc{Sample}(\mathcal{C}_{(k)},1)$ ,  $\mathcal{S}_\tau \leftarrow \mathcal{S}_\tau \cup \{E\}$\;
  }
}
\Else{
  {\color{lightblue}{// Coarse-grained allocation}}\;
  \For{$k\leftarrow 1$ \KwTo $K$}{
    % {\color{lightblue}{// At least one sample per class}}\;
    $a_k\leftarrow 1$, 
    % {\color{lightblue}{// Proportional distribution by class sizes}}\;
    $n_k\leftarrow |\mathcal{C}_k|$, $p_k\leftarrow n_k/\sum_{j=1}^{K} n_j$\;
    $a_k\leftarrow a_k+\lfloor (Q-K)\cdot p_k\rfloor$\;
    \If{$a_k>n_k$}{
      $a_k\leftarrow n_k$\;
    }
  }

  {\color{lightblue}{// Fine-grained adjustment}}\;
  $\sigma \leftarrow \textsc{Argsort}_{\mathrm{desc}}(\{p_k\}_{k=1}^{K})$, 
  $L \leftarrow Q-\sum_{k=1}^{K} a_k$\;
  \While{$L>0$ \textbf{and} $\exists k: a_k<n_k$}{
    \For{$t\leftarrow 1$ \KwTo $K$}{
      $k\leftarrow \sigma(t)$\;
      \If{$a_k<n_k$ \textbf{and} $L>0$}{
        $a_k\leftarrow a_k+1$, $L\leftarrow L-1$\;
      }
    }
  }

  {\color{lightblue}{// In-class sampling}}\;
  $\mathcal{S}_\tau \leftarrow \emptyset$\;
  \For{$k\leftarrow 1$ \KwTo $K$}{
    $\mathcal{S}_\tau \leftarrow \mathcal{S}_\tau \cup \textsc{Sample}(\mathcal{C}_k,a_k)$\;
  }
}
\Return $\mathcal{S}_\tau$\;
\end{algorithm}
}

\subsection{Problem Restatement and Notation}
We consider a budgeted experience selection problem for a fixed tool $\tau$.
Let $\mathcal{E}_\tau=\{E_i\}_{i=1}^{N_\tau}$ be the experience pool, where
$N_\tau \triangleq |\mathcal{E}_\tau|$.
We select $Q$ experiences under a context-length budget. 
In our structured definition, each experience $E$ is associated with a binary score vector
\begin{equation}
\kappa(E)\triangleq scores(E)\in\{0,1\}^d,
\end{equation}
where $d$ is the score dimension.
Experiences are grouped into equivalence classes according to identical score vectors.
Let $\{\mathcal{C}_k\}_{k=1}^{K}$ denote the induced partition of $\mathcal{E}_\tau$, and let $K$ be the number of
observed equivalence classes.
The class distribution on the experience pool $\mathcal{E}_\tau$ is $p_k \triangleq {|\mathcal{C}_k|}/{N_\tau}$ for $k=1,\dots,K$, and $p_k>0$.  
After selection, we obtain a distribution $q=(q_1,\dots,q_K)$ in the selected set over the same $K$ classes,
where $q_k$ denotes the fraction of selected experiences belonging to class $\mathcal{C}_k$ and
$\sum_{k=1}^K q_k = 1$.

Our selection method can be viewed at the class-distribution level as forming a convex combination between
the uniform distribution (encouraging coverage) and the original distribution (preserving representativeness), denoted as
\begin{equation}
\label{eq:q_ours_mix}
q = \alpha u + (1-\alpha)p,
\qquad \alpha \triangleq \frac{K}{Q},
\end{equation}
where $u\in\Delta^{K-1}$ denotes the uniform distribution over the $K$ observed equivalence classes, and
$u_k \triangleq 1/K$, and we focus only on the common regime $Q\ge K$ where full class coverage is feasible. Table~\ref{tab:ecs_notation} summarizes the notations.


\begin{table}[t]
\caption{Notations.}
\centering
\small
\begin{tabular}{@{}lp{0.68\linewidth}@{}}
\toprule
Symbol & Meaning \\
\midrule
$\tau$ & A fixed tool \\
$\mathcal{E}_\tau$ & Experience pool for tool $\tau$ \\
$N_\tau$ & Number of experiences, $N_\tau=|\mathcal{E}_\tau|$ \\
$d$ & Score dimension, $\kappa(E)\in\{0,1\}^d$ \\
$\mathcal{C}_k$ & $k$-th equivalence class \\
$K$ & Number of observed equivalence classes \\
$Q$ & Selection quota \\
$p$ & Class distribution of experience pool, $p_k=|\mathcal{C}_k|/N_\tau$ \\
$q$ & Class distribution of selected experiences\\
$u$ & Uniform distribution over $K$ classes, $u_k=1/K$ \\
$\alpha$ & Mixture weight, $\alpha=K/Q$ \\
\bottomrule
\end{tabular}

\label{tab:ecs_notation}
\end{table}

\subsection{Metrics and Bounds}
The equivalence-class-based selection method aims to induce a distribution $q$ that (1) remains representative of the
original experience pool while (2) covering as many distinct invocation patterns as possible given quota $Q$.
We quantify these two objectives using (i) distribution alignment and (ii) diversity, and derive lower bounds to show the effectiveness of {\ours}.

\subsubsection{Distribution Alignment}
Distribution alignment can be measured by the negative KL divergence~\cite{kl}
\begin{equation}
\text{Align}(q,p)\triangleq -D_{\mathrm{KL}}(q\|p)
= -\sum_{k=1}^{K} q_k \log\frac{q_k}{p_k},
\end{equation}
where larger values indicate better alignment. 

Under Eq.~\eqref{eq:q_ours_mix}, we lower bound $\text{Align}(q,p)$ by upper bounding $D_{\mathrm{KL}}(q\|p)$.
Using the convexity of $D_{\mathrm{KL}}(\cdot\|p)$ in its first argument, we have
\begin{align}
D_{\mathrm{KL}}(q\|p)
&= D_{\mathrm{KL}}\!\left(\alpha u + (1-\alpha)p \,\middle\|\, p\right) \nonumber\\
&\le \alpha D_{\mathrm{KL}}(u\|p) + (1-\alpha)D_{\mathrm{KL}}(p\|p) \nonumber\\
&= \alpha D_{\mathrm{KL}}(u\|p). \label{eq:kl_upper}
\end{align}
Equivalently,
\begin{equation}
\label{eq:align_lower}
\text{Align}(q,p) \ge -\alpha D_{\mathrm{KL}}(u\|p).
\end{equation}

As the quota $Q$ increases, $\alpha=K/Q$ decreases, the bound in Eq.~\eqref{eq:kl_upper} tightens, and $q$ becomes progressively closer to the original distribution $p$. 
In addition, our fine-grained adjustment reallocates any leftover quota to other classes that still have available instances. 
This avoids wasting quota due to over-allocation to small classes and helps keep the induced distribution $q$ closer to the original distribution $p$, which further improves the $\text{Align}(q,p)$. The distribution alignment $\text{Align}(q,p)$ achieves its upper bound of $0$ when the selected distribution $q$ exactly matches the original distribution $p$. 
This design preserves representativeness of the selected experiences by allocating more quota to larger equivalence classes, which correspond to more frequent invocation patterns in the original experience pool of tool $\tau$.

% 因为占比越大的等价类代表了越经常出现的调用范式，这使得选择出的经验具有代表性。
% 我们算法中的Fine-grained adjustment将多余的配额二次发送，消除了随机采样因为类别数目不够带来的空白。使得DISTRIBUTION ALIGNMENT尽可能靠近这个上限，并在q=p的时候达到最大。

\subsubsection{Diversity}
Diversity can be measured using Shannon entropy~\cite{entropy}
\begin{equation}
\text{Div}(q)\triangleq H(q)
= -\sum_{k=1}^{K} q_k \log q_k,
\end{equation}
where larger values correspond to more uniform coverage over invocation patterns. 

Under Eq.~\eqref{eq:q_ours_mix}, we can guarantee a non-trivial level of diversity by lower bounding $H(q)$.
Using the concavity of entropy,
\begin{align}
H(q)
&= H\!\left(\alpha u + (1-\alpha)p\right) \nonumber\\
&\ge \alpha H(u) + (1-\alpha)H(p) \nonumber\\
&= \alpha \log K + (1-\alpha)H(p). \label{eq:entropy_lower}
\end{align}
Therefore, our method guarantees $\text{Div}(q)$ is at least a convex combination of maximal diversity $\log K$ and the original diversity $H(p)$, controlled by $\alpha=K/Q$, and $\text{Div}(q)$ approaches its maximum value $\log K$ when the selection quota $Q$ is close to the number of equivalence classes $K$. This ensures that the selected experiences maintain broad coverage over invocation patterns even under a limited quota, which is particularly beneficial when the original distribution $p$ is long-tailed. In such cases, it mitigates the dominance of a few high-frequency patterns and preserves coverage of rare yet informative ones.

\subsection{Discussion}
% 基于等价类的划分是一种扩展性很强的方法。其能够很自然地利用LLM评估得到的多维向量，将不同方面的质量综合起来，整理出不同的类别范式以供使用。继承自这种划分，基于等价类的选择能够灵活地平衡对常见范式的关注程度和对不同范式的覆盖性。该方法可以扩展至广阔的应用场景中，从而在不改变下游目标的情况下更显式地控制所选子集的代表性和多样性。比如样本选择，和损失定义。
% Equivalence class partitioning is a highly scalable approach. By aggregating quality signals across multiple dimensions, it naturally organizes the results of LLM-as-a-Judge into distinct pattern categories for downstream use. Building on this partitioning, an equivalence-class-based selection method can flexibly balance prioritizing common patterns with covering diverse ones. This method can be applied to a wide range of tasks that require explicit control over the representativeness and diversity of a selected subset, such as sample selection and loss-function design.
As a robust and inherently scalable architecture, equivalence class partitioning redefines how quality signals are aggregated. It systematically maps the outputs of LLM-as-a-Judge into a rigorous taxonomy of pattern categories. This foundational partitioning empowers a superior selection mechanism that masters the critical equilibrium between representativeness and diversity. Beyond mere subset selection, this framework serves as a universal engine for any task requiring precise distributional control, ranging from strategic data mixture to the fundamental optimization of loss functions. 
% 很自然的于工具调用的错误归因联系起来。但该形式本质上是通用的：只要能够通过某个评分函数 $\kappa(\cdot)$ 将样本映射到离散模式（等价类），并且存在固定预算 $Q$（例如上下文长度、存储容量或标注成本），就可以用相同思路进行子集构建。进一步地，在许多下游方法中，训练/优化往往在样本层面定义损失或奖励（loss/reward）并据此进行采样或加权；此时也可以引入本文的等价类视角，将原本``基于样本''的策略扩展为``先在类级别选择/分配，再映射回样本''的实现方式，从而在不改变下游目标的情况下更显式地控制代表性与覆盖度。% 基于将等价类的划分方法能够自然的将LLM评估得到的二元分数与工具调用的模式对应起来，从而支持基于等价类的

\section{Baseline Details}\label{app:baselines}
\paragraph{Few-shot Learning.} A prompting method~\cite{fewshot} that provides a small number of tool-use examples in the prompt to let the model learn from them. To ensure a fair comparison, we use the same invocations in the train set as in {\ours} as examples, and the same truncation for each tool according to the context length.
\paragraph{DRAFT.} Dynamically Refining
tool documentation through the Analysis of Feedback and Trials (DRAFT)~\cite{Draft} is a framework designed to bridge the comprehension gap between LLMs and external tools through an iterative and self-improving methodology. Unlike our approach, DRAFT decouples the process from the agent's inference procedure and improves tool documentation quality through three steps: experience gathering, learning from experience, and documentation rewriting. Since it requires the tool to be actually executed in order to function, on the tools that cannot be actually executed in the MetaTool dataset and parts of the BFCL-V3 dataset, we use the tool invocation information we extracted in {\ours} instead of the real execution. We use Qwen3-max as the agent model, consistent with our setup.
\paragraph{Mem0.} Mem0~\cite{chhikara2025mem0} is a memory layer designed to provide agents with a persistent knowledge base. In our experiments, we specifically utilize its \texttt{procedural\_memory} module (invoked via the \texttt{add} method with \texttt{memory\_type="procedural\_memory"}). As its native design cannot reliably capture tool-level signals (it mainly stores user preferences or coarse-grained agent behaviors), we explicitly provide the interactions between agents and tools from the training set. When retrieving the memory during inference, we apply the same truncation for each tool according to the context length. We use Qwen3-max as the agent model, consistent with our setup.
% By using Mem0 as a baseline, we evaluate the performance of a standard procedural memory storage and retrieval system against \ours, which introduces more complex distillation and utility-based refinement mechanisms.
\label{tried}


\section{Dataset and Metric}\label{app:dataset}
\paragraph{MetaTool} MetaTool~\citep{metatool} is a benchmark designed to evaluate whether LLMs have tool usage awareness and can correctly choose tools. 
We use four sub-task settings that test tool selection to construct a dataset centered on the 100 most frequently used tools, with a maximum of 10 tasks per sub-task involving the same tool, forming a total of 1,669 tasks. Specifically, $\textit{Similar}$ sub-task setting denotes tool selection tasks with similar tool candidates (500 tasks). $\textit{Scenario}$ sub-task setting denotes tool selection tasks in specific scenarios (530 tasks). $\textit{Reliability}$ sub-task setting denotes tool selection tasks with possible reliability issues, where the ground truth tool is excluded from the candidate list (500 tasks). $\textit{Multi-tool}$ sub-task setting denotes multi-tool selection tasks (139 tasks). 
For evaluation, we follow the benchmark's matching method to assess whether the set of tools selected by the agent exactly matches the ground truth answers.


\paragraph{API-Bank} API-Bank~\citep{apibank} is a benchmark designed to evaluate the end-to-end performance of tool-augmented LLMs, specifically in planning, retrieving, and invoking tools within dialogue systems. 
It simulates a comprehensive tool-use environment containing 53 diverse tools and includes 399 tasks in its Level-1 subset, which specifically targets the correctness of tool invocations. 
For evaluation, we follow the dataset's designed matching method, focusing on precision by checking whether the predicted tool name and parameter values exactly match the ground truth annotations.
% In our experiments, we adopt the standard evaluation protocol of API-Bank Lv1 to measure the agent's success in correctly formulating API invocations based on user instructions.


\paragraph{BFCL-V3} The Berkeley Function Calling Leaderboard V3 (BFCL-V3)~\citep{patilberkeley} is a comprehensive benchmark assessing tool invocation capabilities across diverse programming languages and interaction patterns. 
It comprises varied categories, ranging from simple function invocations to complex scenarios involving multiple or parallel invocations. Because some of these tools have only been invoked one time, we aggregate the tasks according to the use of the tools in the ground truth. We then filter the data to keep only tasks in which every tool used also appears in at least two other related tasks, and we randomly sample up to 50 tasks per category, excluding the ``Long Context'' category to avoid confounding effects from different context-length limits. The remaining 461 tasks constitute the test set. To build the training set, we gather all corresponding related tasks and, for each tool, randomly sample 10 tasks when it is associated with more than 10, resulting in 1,542 training tasks. In our analysis, we divide the tasks into four sub-tasks according to (1) the number of inference turns required and (2) whether the task involves problematic settings, such as testing tool-use relevance (irrelevance), missing tools, or missing required information for tool invocation. Specifically, $\textit{Normal}\mathrm{}_{\mathrm{single}}$ sub-task setting denotes single-turn baseline tasks under normal settings (255 tasks). $\textit{Normal}\mathrm{}_{\mathrm{multi}}$ sub-task setting denotes multi-turn baseline tasks under normal settings (50 tasks). $\textit{Problem}\mathrm{}_{\mathrm{single}}$ sub-task setting denotes single-turn tasks with problematic settings (56 tasks). $\textit{Problem}\mathrm{}_{\mathrm{multi}}$ sub-task setting denotes multi-turn tasks with problematic settings (100 tasks). 
For evaluation, the benchmark employs a rigorous dual-evaluation framework: Abstract Syntax Tree matching to verify syntactic correctness and executable testing to ensure the functional outcome matches the ground truth. A task is considered successful only if every round and step passes both Abstract Syntax Tree matching and executable testing.
\section{Implementation details}\label{app:tried}
\subsection{Settings about Experience}
In the main experiments, we use the prompt template for tool invocation to collect experience with the 10-dimensional vector $scores$. For the MetaTool dataset, the answer has only the tool name and no real execution of the tool. Thus, we include a simple prompt ``This is a special tool selection task, you only need to evaluate the tool-use environment, and the agent behavior and the tool reliability should be assumed to be perfect and get all the points by default.'' that guides the agent to focus only on tool selection. For tools in the dataset that can be executed, we use their actual execution results and returned outputs as the tool response field in metadata $m$. This returned content may be vague or even empty, which constitutes one of the challenges that {\ours} is designed to handle. For tools that cannot be executed, we instead use the ground truth as the tool response. When validating the effectiveness of our equivalence-class-based selection method, since the experiences we collected all have $d=10$, we compute the expected probability of 1 for a single dimension, and use a Bernoulli distribution~\cite{Bernoulli} to simulate cases where $d$ takes other values. 


\subsection{Hyperparameter Quota $Q$}
Given a tool $\tau$, we aim to select a safe maximum quota under a fixed context-length budget $B$ (e.g., the maximum tokens allowed for ${\text{LLM}}_{\text{summarizer}}$ beyond the prompt). Let $\mathcal{E}_\tau=\{E_i\}_{i=1}^{N_\tau}$ be the experience pool for $\tau$. We define the length $\ell(E_i)$ of an experience as the token length of its serialized form used at inference time.
We first sort experiences by length in descending order. Let $\pi$ be a permutation such that
\begin{equation}
\small
\label{eq:sort_desc}
\ell\!\left(E_{\pi(1)}\right)\ge \ell\!\left(E_{\pi(2)}\right)\ge \cdots \ge \ell\!\left(E_{\pi(N_\tau)}\right).
\end{equation}
We then compute the target selection quota $Q_\tau$ via a greedy procedure that iterates from the longest experience to the shortest and stops once the budget would be exceeded, as denoted:
\begin{equation}
\small
\label{eq:quota_greedy}
Q_\tau \triangleq \max\left\{q\in\{0,\ldots,N_\tau\} \;\middle|\; \sum_{j=1}^{q}\ell\!\left(E_{\pi(j)}\right)\le B_\tau \right\}.
\end{equation}

As a result, $Q_\tau$ is the safe maximum number of experiences that can be accommodated within budget $B$, which is used for subsequent experience distillation and experience reuse phases.

\subsection{Use of Guidance}
The harness uses the guidance produced by {\ours} in two forms. For one tool at one time, only one form is used. 
Dynamic guidance is the default, since the experience pool can still change. The harness appends it after the tool description as a contextual prompt. The prompt is: ``This is optional guidance on how to better use the tool. You may refer to it selectively: \{dynamic guidance\}.'' The agent may follow it when the current context matches, and a later trajectory can replace it.
Stable guidance is written only after the criterion in Eq.~\ref{eq2} holds, that is, when a consistent non-perfect invocation pattern dominates. The harness stores it with the tool and inserts it into the schema as a constraint: ``This is guidance on using the tool learned from previous experience. You must follow it: \{stable guidance\}.'' This keeps a validated capability boundary or failure pattern in force across tasks. If a new equivalence class later dominates the pool, the harness replaces that schema constraint with the updated guidance.

\begin{table*}[t]
\centering
\caption{Ablation study of alternative variants across different phases of {\ours}. All numbers are Pass@1 on BFCL-V3 with Qwen3-8B. Each row replaces one component of the original design.}
\resizebox{0.9\linewidth}{!}{
\begin{tabular}{llcc}
\toprule
\textbf{Phase} & \textbf{Variant} & \textbf{Accuracy} & \textbf{Reduction vs Original} \\
\midrule
\multirow{1}{*}{Original}
& Full method & \textbf{67.46} & 0.00 \\
\midrule
\multirow{2}{*}{Experience Acquisition}
& Directly use full trajectories (no experiences) & 60.08 & -7.38 \\
& Directly prompt the LLM to extract experiences & 63.34 & -4.12 \\
\midrule
\multirow{4}{*}{Experience Distillation}
& Use only the most recent 1 experience & 62.26 & -5.20 \\
& Keep only the most recent 3 experiences & 64.64 & -2.82 \\
& Remove the experience filter step & 65.94 & -1.52 \\
& Remove the experience summary step & 63.56 & -3.90 \\
\midrule
\multirow{2}{*}{Experience Reuse}
& Only dynamic strategy & 66.37 & -1.09 \\
& Only stable strategy & 63.56 & -3.90 \\
\bottomrule
\end{tabular}
}
\label{tab:variants}
\vspace{-1mm}
\end{table*}




\section{Things We Tried}\label{app:variants}

% 在\{ours}的设计和效果验证时，我们尝试了很多其它方法，这些方法可能没有起作用或不符合我们的设定。这是我们能记住的东西。\textbf{guidance添加位置}。我们尝试将获取到的guidance按工具列表顺序添加到角色为system的提示词最后，这略微降低了我们的效果，我们推测是因为这使得模型增加了一个``从工具列表中去system prompt找到对应工具的guidance''的过程。\textbf{将groundtruth作为采样得到的正向经验加入reme}.这使得我们的方法的效果略微有一点提升（不到1%），说明{\ours}已经可以在没有groundtruth的时候很好的总结工具能力和最佳实践。\textbf{使用更严格的tool response}。我们将数据集evalution的结果加入工具响应，并提示模型``你应该这样完成这个任务...''，这使得我们的方法有很大波动，经过人工排查发现数据集本身质量问题会提供噪声。\textbf{尝试了更多epoch.}我们在实验中的数据集上测试了使用我们的方法进行多epoch推理，每一轮会使得\{ours}得到同一个任务上不同的经验。我们发现{\ours}在第二epoch后几乎不会再有提升，说明reme能在很快时间内收敛，而无法做对的任务则反应了模型和我们方法的上限。\textbf{{使用我们在经验获取阶段获取的经验存入Mem0。}这使得Mem0有了较大进步，能够和我们的方法相比较，但原理上来看这和我们的方法区别不大且并非Mem0的官方实现。\textbf{和基于微调的方法比较。}我们还将我们的方法和两个能找到的基于微调的方法进行了比较，分别是Toolllama和huggingface上发布的qwen3-8b模型，前者因为接口的落后无法与当下任务匹配而取得了很低的分数（几乎是baseline方法的一半），后者尽管相比基础模型有了一定改进（大约3%），但因起没有披露详细的训练方法而被移除。


During the design and evaluation of {\ours}, we explored many alternative variants. Some did not work or did not align with our experimental setting, and here is what we remember:


\textbf{Alternative design variants.}
We explore several alternative variants across different phases of {\ours} to better understand the impact of each design choice. Each row replaces exactly one component of the original design without combining modifications. The results are summarized in Table~\ref{tab:variants}.

\textbf{Guidance insertion position.}
We append the generated guidance to the end of the system prompt following the order of the tool list. This slightly degrades performance. We conjecture that it introduces an extra step for the model to look up the corresponding guidance for each tool from the system prompt based on the tool list.

\textbf{Adding ground truth as positive trajectories into {\ours}.}
Acquiring experiences from ground truth yields a small improvement (less than 1\%). This suggests that {\ours} can already summarize a tool's capability boundaries and best practices well without ground truth or fine-grained feedback.
% {采样标准差} 我们统计了主实验中的三次采样的标准差，尽管在部分设定上有轻微出入，我们的方法在大多数设定上相比其它所有方法都取得了最小的标准差，从另外一个角度验证了{\ours}对agent调用工具鲁棒性的增强。
% \textbf{Using stricter tool responses.}
% We augmented tool responses with dataset evaluation results. We also instructed the model with directives such as ``you should complete this task in this way \ldots''. This led to large performance fluctuations. Manual inspection indicates that quality issues in the dataset can inject substantial noise.

\textbf{Standard deviation.}
We compute the standard deviation over the three runs in our experiment. Our method achieves the smallest standard deviation compared to all other baselines in most settings, providing additional evidence that {\ours} improves the robustness of the agent's tool use.

\textbf{More epochs.}
We test multi-epoch inference on our experimental datasets. Each epoch allows {\ours} to acquire different experiences for the same task. We observe almost no further improvement after the second epoch. This indicates that {\ours} converges quickly, and tasks that remain unsolved reflect the upper bound of the backbone agent and our method.
% \textbf{Storing experiences acquired in the experience acquisition stage into Mem0.}
% This modification substantially improved Mem0, making it comparable to our method. However, we do not report these results because they do not align with mem0's design and, conceptually, they amount to an alternative implementation of our approach.

\begin{table}[t]
\caption{Comparison with tuning-based baselines (Qwen3-8B, \%).}
\vspace{-1mm}
\centering
\fontsize{8}{9}\selectfont
\setlength{\tabcolsep}{1.5pt}
\begin{tabular*}{\columnwidth}{@{\extracolsep{\fill}}lccc@{}}
\toprule
\textbf{Method} & \textbf{MetaTool} & \textbf{API-Bank} & \textbf{BFCL-V3} \\
\midrule
No Method & 76.51 & 83.96 & 58.79 \\
ToolLLM & 38.23 & 42.11 & 24.95 \\
Qwen3-8B-sft & 79.51 & 86.97 & 58.13 \\
{\ours} & 83.52 & 86.47 & 67.46 \\
\bottomrule
\end{tabular*}
\label{tab:tuning}
\vspace{-2mm}
\end{table}

\textbf{Comparison with tuning-based methods.}
We compare {\ours} with two available tuning-based baselines, ToolLLM~\cite{11} and Qwen3-8B-sft~\footnote{\url{https://huggingface.co/1nstaller/qwen3_8b_sft_tool_use}}, on Qwen3-8B. Table~\ref{tab:tuning} reports Pass@1. ToolLLM is below No Method on all three datasets. Its backbone is outdated, and the output format becomes rigid after training. Qwen3-8B-sft is higher than No Method on MetaTool and API-Bank, and slightly lower on BFCL-V3. These methods are not in the main table: tuning shifts the model toward its training corpus, which is a different setting from runtime tool-use robustness.

\textbf{More challenging tasks.}
We also verify the effectiveness of our method on other benchmarks. Under the ``Oracle'' setting of MCPVerse~\cite{lei2025mcpverse} using Qwen3-32B as the backbone model, our method improves the average accuracy from 23.00\% to 26.10\%. However, for smaller models, the gains are substantially constrained by agent capability and context limitations. To more clearly show our method's gains in tool use, we do not choose these challenging tasks as our main experimental benchmarks.
% {更有挑战性的任务。}我们在mcp相benchmark上也验证了我们方法的有效性，在mcpverse的Oracle设定中，使用Qwen3-32B模型作为backbone，我们的方法将平均准确率从23.00%提高到26.10%体现了其有效性。然而在较小尺寸的模型上，该提升大大受到agent能力和上下文的制约，为了验证我们方法的通用性，我们没有选择这些对agent而言难度很大的任务作为实验benchmark。

\begin{table}[t]
\caption{Ablation of the prompt (Qwen3-8B, \%).}
\vspace{-1mm}
\centering
\fontsize{8}{9}\selectfont
\setlength{\tabcolsep}{2pt}
\begin{tabular*}{\columnwidth}{@{\extracolsep{\fill}}lcc@{}}
\toprule
\textbf{Prompt} & \textbf{API-Bank} & \textbf{BFCL-V3} \\
& \multicolumn{2}{c}{Pass@1 / Avg@3 / Pass@3} \\
\midrule
No Method & 83.96/83.88/84.21 & 58.79/58.28/60.30 \\
w/o tool-use environment & 85.21/85.38/85.46 & 65.29/63.12/67.03 \\
w/o agent behavior & 83.96/84.13/84.21 & 64.43/61.97/65.29 \\
w/o tool reliability & 85.71/85.88/85.96 & 65.08/62.55/65.73 \\
Full prompt & \textbf{86.47/87.89/87.97} & \textbf{67.46/66.96/68.33} \\
\bottomrule
\end{tabular*}
\label{tab:score_groups}
\vspace{-2mm}
\end{table}

\textbf{Prompt ablation.}
We ablate the evaluator prompt on Qwen3-8B by removing one score group at a time, and Table~\ref{tab:score_groups} reports the results. Removing any group lowers the scores 
relative to the full prompt, and removing agent behavior causes the largest drop on 
both datasets. MetaTool only scores tool selection, so that dataset already uses the 
tool-use environment items alone.


\begin{table*}[t]
\centering
\caption{Equivalence classes of \texttt{BookHotel}.}
\label{tab:bookhotel}
\small
\begin{tabular}{clp{0.72\linewidth}}
\toprule
\textbf{Count} & \textbf{Class} & \textbf{What the scores show} \\
\midrule
6 & [1,1,1,1,1,1,1,1,1,1] & The scene and the parameters are correct, and the booking succeeds. The user confirms a hotel booking with valid dates and guest counts, and the tool returns an order id. \\
1 & [0,1,1,1,1,1,1,1,1,1] & Tool-use environment: the call succeeds, but the timing is wrong. The user is still asking for prices and has not confirmed, yet \texttt{BookHotel} is called. \\
1 & [1,1,1,1,1,0,1,1,1,1] & Agent behavior: the call succeeds, but a parameter drifts from the context. The user asks for Shanghai Hilton, and only Hilton is filled. \\
1 & [1,1,1,0,1,0,0,0,0,0] & Agent behavior and tool reliability: an illegal parameter makes the call fail. \texttt{room\_count} is out of range, and the tool returns an error. \\
3 & [1,1,1,1,0,1,0,0,0,0] & Agent behavior and tool reliability: the meaning is right but the format is wrong. \texttt{adult\_count} is written as a string, and the call fails. \\
1 & [0,0,1,1,1,0,0,0,0,0] & All three groups: the scene, the tool, and the parameter are wrong. The user wants a meeting room, but the agent calls \texttt{BookHotel} with \texttt{hotel\_name} set to Conference Room, and the call fails. \\
\bottomrule
\end{tabular}
\end{table*}


% \newpage
\section{Case Analysis}\label{case}

\subsection{Equivalence Classes of One Tool}
\texttt{BookHotel} takes the number of guests and the dates, and returns an order id when the booking succeeds. Table~\ref{tab:bookhotel} lists the equivalence classes in its experience pool. The score vector follows the evaluator order: tool-use environment (2), agent behavior (4), and tool reliability (4).

\subsection{{\ours} Improves the Agent's Tool-Use Ability from Multiple Aspects}
We summarize representative cases of tool-use improvements brought by our method into six aspects. For each aspect, we provide an illustrative example with a brief explanation.
\paragraph{(1) Improve Tool-Use Awareness.}

When the user implicitly asks about something related to a historical period, the agent tends to answer directly from its internal knowledge in a conversational manner instead of invoking the \texttt{timeport} tool. {\ours} learns from past experience that the tool should be invoked in such cases.

\paragraph{(2) Distinguish Similar Tools.}
The agent uses \texttt{RestaurantBookingTool} to make restaurant reservations. When the user's request is ambiguous, the agent often confuses it with the semantically similar \texttt{local} tool, since \texttt{local} can be used to search for restaurants. {\ours} learns the capability boundary from past experience, i.e., it should invoke the booking tool only when the user explicitly asks to reserve a table.

\paragraph{(3) Invoke Tools Correctly.}
The agent uses the \texttt{inventory\_management} tool to complete inventory management tasks. {\ours} learns from past experience the constraints on the allowed parameter content and required formats for this tool.

\paragraph{(4) Fix Tool-Use Errors.}
The agent uses \texttt{EmergencyKnowledge} to retrieve possible diseases, but the first invocation fails because the agent provides an overly descriptive input. After {\ours} recognizes this issue, it fixes the input formulation in subsequent invocations.

\paragraph{(5) Identify Tool Dependencies.}
The agent uses the file-system tool \texttt{wc} to count the number of lines in a file. {\ours} learns from past experience that the error ``file does not exist'' occurs frequently. The guidance encourages the agent to run \texttt{ls} or \texttt{find} beforehand, which can effectively avoid such failures. 

\paragraph{(6) Understand Tool Reliability.}
From prior experiences where the agent uses \texttt{QueryRegistration} to check patient appointment registrations, {\ours} observes frequent failures caused by tool availability issues and identifies them as tool-side problems. It therefore reminds the agent to check tool availability before invocation or to seek alternative tools. 
% In a later query, the agent completes the task using an alternative tool.

\subsection{Cases Where {\ours} Fails}

\paragraph{(1) Unseen Situations.}
If the agent uses an incorrect tool, or fails to use a tool when it should, {\ours} cannot collect trajectories of the correct tool invocation and therefore cannot learn from it to improve. As a result, there is no experience or guidance for that tool. For previously seen tools, {\ours} may still fail when experience is limited: some invocation patterns or error modes may not have appeared in the past trajectories, leaving the agent without relevant guidance to prevent them.
% \begin{figure}[!h]
%     \centering
%     \includegraphics[width=0.94\linewidth]{figures/badcase1.png}
%     \caption{Case of failed tool use: Unseen Tools.}
% \end{figure}

\paragraph{(2) Hard-to-analyze Trajectories.}
% A user asks the agent to set an account password. After multiple rounds of tool interaction, the agent finds that no account exists and decides to register one. This logic may be reasonable for many agents, so {\ours} also fails to summarize that the user intent is merely to initialize the account. If no account exists, the agent should further communicate with the user. Such issues are stable to capture because they stem from the agent's ability to precisely infer user expectations in every step.
% agent调用edit_ticket完成用户``帮我编辑3号票据内容为出货信息。''的任务。然而智能体判断的出货信息包括所有的货物名称，数量和价格明细等信息，与用户所期望的符合模板等物流信息不同，所以即使该工具被成功调用，用户的要求仍然没有完成。该错误是由于智能体对用户意图的判断失误而不是工具调用错误而引起的，属于{\ours} 可以解决范围之外的问题。
The agent invokes \texttt{edit\_ticket} to handle ``Please edit ticket No. 3 to contain shipping information.'' It fills the ticket with all item names, quantities, and price{\parfillskip=0pt\par} 

\begin{figure*}[p]
    \centering
    \includegraphics[width=0.94\linewidth]{figures/case1.png}
    \caption{Successful case of improved tool use: improve tool-use awareness.}
    \vspace{-2mm}
\end{figure*}

\begin{figure*}[p]
    \centering
    \includegraphics[width=0.94\linewidth]{figures/case2.png}
    \caption{Successful case of improved tool use: distinguish similar tools.}
    \vspace{-2mm}
\end{figure*}

\begin{figure*}[p]
    \centering
    \includegraphics[width=0.94\linewidth]{figures/case3.png}
    \caption{Successful case of improved tool use: invoke tools correctly.}
    \vspace{-2mm}
\end{figure*}

\begin{figure*}[p]
    \centering
    \includegraphics[width=0.94\linewidth]{figures/case4.png}
    \caption{Successful case of improved tool use: fix tool-use errors.}
    \vspace{-2mm}
\end{figure*}

\begin{figure*}[p]
    \centering
    \includegraphics[width=0.94\linewidth]{figures/case5.png}
    \caption{Successful case of improved tool use: identify tool dependencies.}
    \vspace{-2mm}
\end{figure*}

\begin{figure*}[p]
    \centering
    \includegraphics[width=0.94\linewidth]{figures/case6.png}
    \caption{Successful case of improved tool use: understand tool reliability.}
    \vspace{-2mm}
\end{figure*}

\begin{figure*}[p]
    \centering
    \includegraphics[width=0.94\linewidth]{figures/badcase2.png}
    \caption{Case of failed tool use: hard-to-analyze trajectories.}
    \vspace{-2mm}
\end{figure*}

\begin{figure*}[p]
    \centering
    \includegraphics[width=0.94\linewidth]{figures/badcase3.png}
    \caption{Case of failed tool use: incorrect attributions.}
    \label{fig:badcase}
    \vspace{-2mm}
\end{figure*}

\newpage
\begin{strip}
    \centering
    \includegraphics[width=0.94\linewidth]{figures/badcase4.png}
    \captionof{figure}{Case of failed tool use: over-reaction.}
    \vspace{-2mm}
\end{strip}



\noindent
breakdowns. The user expected that information to be formatted according to a template. The invocation succeeds, but the request is not satisfied. {\ours} learns from the invocation context, arguments, and tool response. These signals look successful, so {\ours} receives no evidence of the mismatch and cannot address it.

\paragraph{(3) Incorrect Attributions.}
% The agent uses \texttt{CancelTimedSwitch} to adjust device status. Because the tool's actual behavior is inconsistent with its schema, the agent never succeeds in invoking it. However, the tool always returns parameter-related errors. During summarization, {\ours} identifies it as a tool-side problem, but in the next attempt it incorrectly attributes the failure to the agent's inappropriate behavior rather than the tool's own issue.
The agent adjusts the device state using \texttt{CancelTimedSwitch}. Because the actual behavior of the tool is not consistent with its schema (the schema specifies that the device id should be passed as the content of parameter ``device\_name''), the agent will never successfully invoke it. However, the tool always returns a ``device not present'' response because it can't find the device. {\ours} incorrectly attributes the failure to inappropriate behavior of the agent, as the agent did not pass in the correct device name, rather than a problem with the consistency of the tool itself.

\paragraph{(4) Over-Reaction.}
When doctors' names always included a prefix, the agent was able to use \texttt{ModifyRegistration} successfully in past experience. Consequently, {\ours} inferred a piece of guidance: always use the name with a title. However, those earlier successes occurred only because the doctor names stored in the system already contained the prefix. When a doctor is entered with a full name rather than a titled name, the agent, following the guidance, still forcibly adds the ``Dr.'' prefix, which leads to a failure.


\subsection{Preventing LLM-Introduced Errors}
Besides the typical failures above, {\ours} uses an LLM, so the LLM can introduce errors of its own. The following defenses address these errors.


\paragraph{Experience distillation.}
The summarizer draws on multiple experiences. A minority of incorrect attributions is corrected by the majority of reliable ones. When the two sides are roughly tied, the summarizer reconsiders the evidence before writing the guidance. In later epochs, experience filtering can partially remove guidance judged unreliable. Human inspection of the guidance-generation process finds that over 95\% of the acquired experiences are reliable, and 98.7\% of the resulting guidance is judged reliable. The remaining guidance, about 1.3\%, may contain an occasional error.


\paragraph{Experience reuse.}
Guidance is optional and dynamic by default. It becomes a schema constraint only when Eq.~\ref{eq2} holds. An occasional mistake stays out of the schema unless the evaluator repeats that same attribution. If an incorrect attribution does become stable guidance, later calls that follow the harmful guidance write the new failure pattern back into the pool. When that pattern becomes the dominant equivalence class, it replaces the previous schema constraint. 


\subsection{Case Coverage Comparison}
% 以API-bank为例，我们选了其中一次推理结果，对{\ours}和其它方法进行了Case-coverage Comparison。图~\ref{overlap}展示了它们相比没有方法的baseline做对和做错样本的数量。在比baseline表现好的部分，我们的方法覆盖了大部分的Few-shot和DRAFT能做对的样本，但和mem0这种对agent进行增强的方法有互补优势。在比baseline Underperform的部分，明显地我们的方法相比其它方法引入了更少的错误，归功于我们Experience Distillation phase对经验进行了effectiveness check并剔除unhelpful的经验。
\begin{figure*}[t]
    \centering
    \includegraphics[width=0.8\linewidth]{figures/overlap.png}
    \caption{Case coverage comparison with baseline methods.}
    \label{overlap}
    \vspace{-2mm}
\end{figure*}
Using API-Bank as the dataset, we select one run as a representative result and compare case coverage between {\ours} and other methods. Figure~\ref{overlap} shows the numbers of cases that each method gets right or wrong compared to the ``No Method'' baseline. 
In the subset where methods outperform the baseline, {\ours} covers most of the cases solved by Few-shot and DRAFT, while exhibiting complementary gains to Mem0, possibly because Mem0 enhances the agent from a process perspective rather than focusing on tool invocation. In the subset where methods underperform the baseline, {\ours} introduces fewer error cases than the other methods. We attribute this robustness to the experience distillation phase, where {\ours} performs an effectiveness check and filters out unhelpful experiences.
%%%%%%%%%%%%%%%%%%%%%%%%%%%%%%%%%%%%%%%%%%%%%%%%%%%%%%%%%%%%%%%%%%%%%%%%%%%%%%%
%%%%%%%%%%%%%%%%%%%%%%%%%%%%%%%%%%%%%%%%%%%%%%%%%%%%%%%%%%%%%%%%%%%%%%%%%%%%%%%


\begin{table}[t]
\caption{Overhead of {\ours} on the three main benchmarks (Qwen3-8B, parallelism $=4$).}
\vspace{-1mm}
\centering
\fontsize{8}{9}\selectfont
\setlength{\tabcolsep}{1.5pt}
\begin{tabular*}{\columnwidth}{@{\extracolsep{\fill}}lcccc@{}}
\toprule
\textbf{Phase} & \begin{tabular}[c]{@{}c@{}}\textbf{Extra}\\\textbf{LLM calls}\end{tabular} & \textbf{Tokens} & \textbf{Cost (CNY)} & \textbf{Time} \\
\midrule
Acquisition & 3,829 & $\approx 6.8\times10^{6}$ & 20.7 & 85\,min \\
Distillation & 3,335 & $\approx 7.2\times10^{6}$ & 24.9 & 74\,min \\
\midrule
Total & 7,164 & $\approx 1.4\times10^{7}$ & 45.6 & 159\,min \\
\bottomrule
\end{tabular*}
\label{tab:cost}
\vspace{-2mm}
\end{table}

\section{Cost Analysis}\label{app:compute}
\paragraph{Full reproduction.}
Full reproduction requires valid API keys for the commercial LLM endpoints referenced in the Experiments section. All backbone, evaluator, and summarizer inference is obtained through these APIs. A CPU-only machine is sufficient for orchestration and local preprocessing. 
In our setup, the driver process, dataset input and output, logging, and the vector database that stores experiences typically remained below approximately 16\,GB of RAM at peak usage. With reduced caching or smaller concurrent batches, 8\,GB of RAM is often sufficient. However, the risk of out-of-memory errors increases if many trajectories are materialized simultaneously. 
Wall-clock time is primarily determined by provider latency, API rate limits, and the specific models or ablations that are enabled. In our environment, running the full set of main experiments required approximately 12 days from start to completion. This duration may vary substantially across different environments and API service tiers. 
Reproducing all tables and baselines reported in the paper is expected to require more than $8\times10^{5}$ remote API or model calls in total.

\paragraph{Online overhead of {\ours}.}
In online use, the experience pool is a vector database indexed by tool name.
On average, each experience takes about 5.7\,KB and each tool about 44\,KB, while a retrieve and an update take about 47\,ms and 170\,ms, respectively.
The extra LLM overhead mainly comes from two calls by the evaluator and the summarizer, with an average extra prompt length of 691 words.
Table~\ref{tab:cost} reports the measured overhead: each tool invocation incurs only an extra \$0.0018 and about 2.5 seconds.
On a sample of 300 tasks, 100 drawn from each dataset and one invocation per tool, Table~\ref{tab:overhead_cmp} compares this overhead with Few-shot, DRAFT, and Mem0. The extra prompt length, extra calls, and extra time of {\ours} stay in the same range as these baselines.

\begin{table}[t]
\caption{Extra overhead on 300 tasks (100 from each dataset, one invocation per tool).}
\vspace{-1mm}
\centering
\fontsize{8}{9}\selectfont
\setlength{\tabcolsep}{1.5pt}
\begin{tabular*}{\columnwidth}{@{\extracolsep{\fill}}lccc@{}}
\toprule
\textbf{Method} & \begin{tabular}[c]{@{}c@{}}\textbf{Prompt length}\\\textbf{(words)}\end{tabular} & \begin{tabular}[c]{@{}c@{}}\textbf{Extra}\\\textbf{LLM calls}\end{tabular} & \begin{tabular}[c]{@{}c@{}}\textbf{Extra}\\\textbf{time (s)}\end{tabular} \\
\midrule
Few-shot & 524 & 0 & 1.3 \\
DRAFT & 629 & 5 & 5.7 \\
Mem0 & 926 & 3 & 8.6 \\
{\ours} & 691 & 2 & 3.1 \\
\bottomrule
\end{tabular*}
\label{tab:overhead_cmp}
\end{table}

\paragraph{Tighter selection quota.}
We also tighten the context budget on BFCL-V3 with Qwen3-8B. Table~\ref{tab:quota_budget} shows that the drop is small: tightening the quota from $Q$ to $Q/8$ lowers Pass@1 from 67.46 to 63.12. A smaller quota makes the $Q/2$ evidence threshold easier to reach. For tools with many equivalence classes, it also makes the same class less likely to recur, which keeps the move from dynamic guidance to stable guidance in check.


\begin{table}[t]
\centering
\caption{\small Pass@1 on BFCL-V3 (Qwen3-8B, \%) when the selection quota is tightened. $\triangle$ is the gain over No Method.}
\label{tab:quota_budget}
\vspace{-1mm}
\fontsize{8}{9.5}\selectfont
\setlength{\tabcolsep}{3pt}
\begin{tabular}{@{}lcccc@{}}
\toprule
\textbf{Quota} & $Q$ & $Q/2$ & $Q/4$ & $Q/8$ \\
\midrule
\textbf{Pass@1} & 67.46 & 66.81 & 65.29 & 63.12 \\
$\triangle$ & 8.67 & 8.02 & 6.50 & 4.33 \\
\bottomrule
\end{tabular}
\vspace{-2mm}
\end{table}

In resource-constrained settings, we recommend using {\ours} in a ``generate once, reuse many times'' manner: because the guidance is synthesized from multiple past experiences, it tends to be broadly applicable, making it possible to avoid these extra costs while maintaining comparable performance.

% \section{Limitations of Our Work}
% % As shown in Table~\ref{tab:summarizer}, replacing an advanced LLM with a smaller one leads to performance degradation.
% While {\ours} substantially improves the robustness of agents' tool use, several limitations remain. 
% First, generating high-quality guidance depends on the capabilities of the evaluator and summarizer agents. This dependency incurs additional computational overhead and may hinder deployment in resource-constrained environments, since each tool invocation triggers two extra agent inferences as well as storage and retrieval costs. In such cases, we recommend using {\ours} in a ``generate once, reuse many times'' manner: because the guidance is synthesized from multiple past experiences, it tends to be broadly applicable, making it possible to avoid these extra costs while maintaining comparable performance. 
% Second, the experience acquisition process depends on the accuracy of the attribution. In complex scenarios where tool feedback is more ambiguous or misleading as shown in Figure~\ref{fig:badcase}, the evaluator may be unable to identify the cause of failure, resulting in incorrect guidance. While we have verified that this issue can be mitigated by aggregating diverse experiences and filtering out unhelpful ones in the second phase, the risk of reinforcing suboptimal behaviors due to incorrect attributions remains. 
% Finally, {\ours} faces a cold-start challenge. The system requires initial interaction trajectories to acquire experiences. If a tool is never invoked, or if certain failure/success patterns are not encountered early on, {\ours} may struggle to acquire the corresponding experiences needed to form effective guidance.

%%%%%%%%%%%%%%%%%%%%%%%%%%%%%%%%%%%%%%%%%%%%%%%%%%%%%%%%%%%%

% \newpage
% \section*{NeurIPS Paper Checklist}


\begin{enumerate}

\item {\bf Claims}
    \item[] Question: Do the main claims made in the abstract and introduction accurately reflect the paper's contributions and scope?
    \item[] Answer: \answerYes{}
    \item[] Justification: The abstract and introduction describe {\texttt{ExpG}}'s three phases (experience acquisition, distillation, reuse), the problem setting (dynamic environments and noisy tool feedback), and empirical claims; the methodology, experiments, analysis, and conclusion parts of the paper match these claims.
    \item[] Guidelines:
    \begin{itemize}
        \item The answer \answerNA{} means that the abstract and introduction do not include the claims made in the paper.
        \item The abstract and/or introduction should clearly state the claims made, including the contributions made in the paper and important assumptions and limitations. A \answerNo{} or \answerNA{} answer to this question will not be perceived well by the reviewers. 
        \item The claims made should match theoretical and experimental results, and reflect how much the results can be expected to generalize to other settings. 
        \item It is fine to include aspirational goals as motivation as long as it is clear that these goals are not attained by the paper. 
    \end{itemize}

\item {\bf Limitations}
    \item[] Question: Does the paper discuss the limitations of the work performed by the authors?
    \item[] Answer: \answerYes{}
    \item[] Justification: An appendix titled ``Limitations of Our Work'' discusses evaluator and summarizer dependency and cost, attribution errors, and cold-start; the ``More Analysis'' subsection (case analysis) further documents failure modes (unseen situations, incorrect attributions, over-reaction, etc.).
    \item[] Guidelines:
    \begin{itemize}
        \item The answer \answerNA{} means that the paper has no limitation while the answer \answerNo{} means that the paper has limitations, but those are not discussed in the paper. 
        \item The authors are encouraged to create a separate ``Limitations'' section in their paper.
        \item The paper should point out any strong assumptions and how robust the results are to violations of these assumptions (e.g., independence assumptions, noiseless settings, model well-specification, asymptotic approximations only holding locally). The authors should reflect on how these assumptions might be violated in practice and what the implications would be.
        \item The authors should reflect on the scope of the claims made, e.g., if the approach was only tested on a few datasets or with a few runs. In general, empirical results often depend on implicit assumptions, which should be articulated.
        \item The authors should reflect on the factors that influence the performance of the approach. For example, a facial recognition algorithm may perform poorly when image resolution is low or images are taken in low lighting. Or a speech-to-text system might not be used reliably to provide closed captions for online lectures because it fails to handle technical jargon.
        \item The authors should discuss the computational efficiency of the proposed algorithms and how they scale with dataset size.
        \item If applicable, the authors should discuss possible limitations of their approach to address problems of privacy and fairness.
        \item While the authors might fear that complete honesty about limitations might be used by reviewers as grounds for rejection, a worse outcome might be that reviewers discover limitations that aren't acknowledged in the paper. The authors should use their best judgment and recognize that individual actions in favor of transparency play an important role in developing norms that preserve the integrity of the community. Reviewers will be specifically instructed to not penalize honesty concerning limitations.
    \end{itemize}

\item {\bf Theory assumptions and proofs}
    \item[] Question: For each theoretical result, does the paper provide the full set of assumptions and a complete (and correct) proof?
    \item[] Answer: \answerYes{}
    \item[] Justification: The equivalence-class-based selection is formalized with notation and objectives in appendix~\ref{proof}, including assumptions on the score partition and budget $Q$, and proofs or sketches for distribution alignment and diversity bounds; the main text references this appendix and states the construction (e.g., Eq.~(\ref{eq:q_ours_mix}) in the appendix).
    \item[] Guidelines:
    \begin{itemize}
        \item The answer \answerNA{} means that the paper does not include theoretical results. 
        \item All the theorems, formulas, and proofs in the paper should be numbered and cross-referenced.
        \item All assumptions should be clearly stated or referenced in the statement of any theorems.
        \item The proofs can either appear in the main paper or the supplemental material, but if they appear in the supplemental material, the authors are encouraged to provide a short proof sketch to provide intuition. 
        \item Inversely, any informal proof provided in the core of the paper should be complemented by formal proofs provided in appendix or supplemental material.
        \item Theorems and Lemmas that the proof relies upon should be properly referenced. 
    \end{itemize}

    \item {\bf Experimental result reproducibility}
    \item[] Question: Does the paper fully disclose all the information needed to reproduce the main experimental results of the paper to the extent that it affects the main claims and/or conclusions of the paper (regardless of whether the code and data are provided or not)?
    \item[] Answer: \answerYes{}
    \item[] Justification: The experiments section specifies benchmarks (MetaTool, API-Bank, BFCL-V3), splits and cross-validation, metrics (\texttt{Avg@3}, \texttt{Pass@3}), backbone models, temperature, $\alpha$, and the roles of $\mathrm{LLM}_{\mathrm{evaluator}}$ and $\mathrm{LLM}_{\mathrm{summarizer}}$; appendices~\ref{app:dataset}, \ref{app:baselines}, and \ref{app:tried} add dataset, baseline, and implementation detail; code and reproduction materials are at \url{https://anonymous.4open.science/r/ExpG-F72B}.
    \item[] Guidelines:
    \begin{itemize}
        \item The answer \answerNA{} means that the paper does not include experiments.
        \item If the paper includes experiments, a \answerNo{} answer to this question will not be perceived well by the reviewers: Making the paper reproducible is important, regardless of whether the code and data are provided or not.
        \item If the contribution is a dataset and\slash or model, the authors should describe the steps taken to make their results reproducible or verifiable. 
        \item Depending on the contribution, reproducibility can be accomplished in various ways. For example, if the contribution is a novel architecture, describing the architecture fully might suffice, or if the contribution is a specific model and empirical evaluation, it may be necessary to either make it possible for others to replicate the model with the same dataset, or provide access to the model. In general. releasing code and data is often one good way to accomplish this, but reproducibility can also be provided via detailed instructions for how to replicate the results, access to a hosted model (e.g., in the case of a large language model), releasing of a model checkpoint, or other means that are appropriate to the research performed.
        \item While NeurIPS does not require releasing code, the conference does require all submissions to provide some reasonable avenue for reproducibility, which may depend on the nature of the contribution. For example
        \begin{enumerate}
            \item If the contribution is primarily a new algorithm, the paper should make it clear how to reproduce that algorithm.
            \item If the contribution is primarily a new model architecture, the paper should describe the architecture clearly and fully.
            \item If the contribution is a new model (e.g., a large language model), then there should either be a way to access this model for reproducing the results or a way to reproduce the model (e.g., with an open-source dataset or instructions for how to construct the dataset).
            \item We recognize that reproducibility may be tricky in some cases, in which case authors are welcome to describe the particular way they provide for reproducibility. In the case of closed-source models, it may be that access to the model is limited in some way (e.g., to registered users), but it should be possible for other researchers to have some path to reproducing or verifying the results.
        \end{enumerate}
    \end{itemize}


\item {\bf Open access to data and code}
    \item[] Question: Does the paper provide open access to the data and code, with sufficient instructions to faithfully reproduce the main experimental results, as described in supplemental material?
    \item[] Answer: \answerYes{}
    \item[] Justification: Code, experiments, and results are shared at \url{https://anonymous.4open.science/r/ExpG-F72B} (also linked in the abstract); appendix~\ref{app:tried} documents key hyperparameters, experience formatting, quota $Q_\tau$, and guidance injection templates, and the repository README documents environment setup and reproduction commands.
    \item[] Guidelines:
    \begin{itemize}
        \item The answer \answerNA{} means that paper does not include experiments requiring code.
        \item Please see the NeurIPS code and data submission guidelines (\url{https://neurips.cc/public/guides/CodeSubmissionPolicy}) for more details.
        \item While we encourage the release of code and data, we understand that this might not be possible, so \answerNo{} is an acceptable answer. Papers cannot be rejected simply for not including code, unless this is central to the contribution (e.g., for a new open-source benchmark).
        \item The instructions should contain the exact command and environment needed to run to reproduce the results. See the NeurIPS code and data submission guidelines (\url{https://neurips.cc/public/guides/CodeSubmissionPolicy}) for more details.
        \item The authors should provide instructions on data access and preparation, including how to access the raw data, preprocessed data, intermediate data, and generated data, etc.
        \item The authors should provide scripts to reproduce all experimental results for the new proposed method and baselines. If only a subset of experiments are reproducible, they should state which ones are omitted from the script and why.
        \item At submission time, to preserve anonymity, the authors should release anonymized versions (if applicable).
        \item Providing as much information as possible in supplemental material (appended to the paper) is recommended, but including URLs to data and code is permitted.
    \end{itemize}


\item {\bf Experimental setting/details}
    \item[] Question: Does the paper specify all the training and test details (e.g., data splits, hyperparameters, how they were chosen, type of optimizer) necessary to understand the results?
    \item[] Answer: \answerYes{}
    \item[] Justification: Training and test construction and evaluation protocols are in the experiments section and appendix~\ref{app:dataset}; models, decoding settings (temperature $0.5$), $\alpha=0.8$, sequential execution, and fixed settings for ablations are in the main text with further detail in appendix~\ref{app:tried}. The method is tuning-free (inference-time guidance), so no optimizer is used for the proposed mechanism.
    \item[] Guidelines:
    \begin{itemize}
        \item The answer \answerNA{} means that the paper does not include experiments.
        \item The experimental setting should be presented in the core of the paper to a level of detail that is necessary to appreciate the results and make sense of them.
        \item The full details can be provided either with the code, in appendix, or as supplemental material.
    \end{itemize}

\item {\bf Experiment statistical significance}
    \item[] Question: Does the paper report error bars suitably and correctly defined or other appropriate information about the statistical significance of the experiments?
    \item[] Answer: \answerYes{}
    \item[] Justification: The main metrics aggregate three independent runs (\texttt{Avg@3}, \texttt{Pass@3}), explicitly quantifying run-to-run variability; appendix~\ref{app:tried} additionally reports standard deviation across those runs as robustness evidence rather than classical significance tests on every table.
    \item[] Guidelines:
    \begin{itemize}
        \item The answer \answerNA{} means that the paper does not include experiments.
        \item The authors should answer \answerYes{} if the results are accompanied by error bars, confidence intervals, or statistical significance tests, at least for the experiments that support the main claims of the paper.
        \item The factors of variability that the error bars are capturing should be clearly stated (for example, train/test split, initialization, random drawing of some parameter, or overall run with given experimental conditions).
        \item The method for calculating the error bars should be explained (closed form formula, call to a library function, bootstrap, etc.)
        \item The assumptions made should be given (e.g., Normally distributed errors).
        \item It should be clear whether the error bar is the standard deviation or the standard error of the mean.
        \item It is OK to report 1-sigma error bars, but one should state it. The authors should preferably report a 2-sigma error bar than state that they have a 96\% CI, if the hypothesis of Normality of errors is not verified.
        \item For asymmetric distributions, the authors should be careful not to show in tables or figures symmetric error bars that would yield results that are out of range (e.g., negative error rates).
        \item If error bars are reported in tables or plots, the authors should explain in the text how they were calculated and reference the corresponding figures or tables in the text.
    \end{itemize}

\item {\bf Experiments compute resources}
    \item[] Question: For each experiment, does the paper provide sufficient information on the computer resources (type of compute workers, memory, time of execution) needed to reproduce the experiments?
    \item[] Answer: \answerYes{}
    \item[] Justification: Appendix~\ref{app:compute} records that runs use commercial LLM APIs with user-supplied keys, that a CPU-only orchestration machine suffices (about 8--16~GB RAM typical in our runs), that full reruns likely exceed $8\times 10^{5}$ API calls, and that our wall-clock was on the order of ten days with high variance across setups; appendix~\ref{app:tried} and \url{https://anonymous.4open.science/r/ExpG-F72B} add implementation detail.
    \item[] Guidelines:
    \begin{itemize}
        \item The answer \answerNA{} means that the paper does not include experiments.
        \item The paper should indicate the type of compute workers CPU or GPU, internal cluster, or cloud provider, including relevant memory and storage.
        \item The paper should provide the amount of compute required for each of the individual experimental runs as well as estimate the total compute. 
        \item The paper should disclose whether the full research project required more compute than the experiments reported in the paper (e.g., preliminary or failed experiments that didn't make it into the paper). 
    \end{itemize}
    
\item {\bf Code of ethics}
    \item[] Question: Does the research conducted in the paper conform, in every respect, with the NeurIPS Code of Ethics \url{https://neurips.cc/public/EthicsGuidelines}?
    \item[] Answer: \answerYes{}
    \item[] Justification: The work uses standard public benchmarks and API-accessible models, follows double-blind conventions, and links artifacts at \url{https://anonymous.4open.science/r/ExpG-F72B}; it does not involve human subjects or sensitive private data collection as part of the contribution.
    \item[] Guidelines:
    \begin{itemize}
        \item The answer \answerNA{} means that the authors have not reviewed the NeurIPS Code of Ethics.
        \item If the authors answer \answerNo{}, they should explain the special circumstances that require a deviation from the Code of Ethics.
        \item The authors should make sure to preserve anonymity (e.g., if there is a special consideration due to laws or regulations in their jurisdiction).
    \end{itemize}


\item {\bf Broader impacts}
    \item[] Question: Does the paper discuss both potential positive societal impacts and negative societal impacts of the work performed?
    \item[] Answer: \answerNA{}
    \item[] Justification: The submission is scoped as a methods and public-benchmark evaluation study and does not include a separate discussion of positive and negative societal deployment impacts; we therefore mark this item not applicable.
    \item[] Guidelines:
    \begin{itemize}
        \item The answer \answerNA{} means that there is no societal impact of the work performed.
        \item If the authors answer \answerNA{} or \answerNo{}, they should explain why their work has no societal impact or why the paper does not address societal impact.
        \item Examples of negative societal impacts include potential malicious or unintended uses (e.g., disinformation, generating fake profiles, surveillance), fairness considerations (e.g., deployment of technologies that could make decisions that unfairly impact specific groups), privacy considerations, and security considerations.
        \item The conference expects that many papers will be foundational research and not tied to particular applications, let alone deployments. However, if there is a direct path to any negative applications, the authors should point it out. For example, it is legitimate to point out that an improvement in the quality of generative models could be used to generate Deepfakes for disinformation. On the other hand, it is not needed to point out that a generic algorithm for optimizing neural networks could enable people to train models that generate Deepfakes faster.
        \item The authors should consider possible harms that could arise when the technology is being used as intended and functioning correctly, harms that could arise when the technology is being used as intended but gives incorrect results, and harms following from (intentional or unintentional) misuse of the technology.
        \item If there are negative societal impacts, the authors could also discuss possible mitigation strategies (e.g., gated release of models, providing defenses in addition to attacks, mechanisms for monitoring misuse, mechanisms to monitor how a system learns from feedback over time, improving the efficiency and accessibility of ML).
    \end{itemize}
    
\item {\bf Safeguards}
    \item[] Question: Does the paper describe safeguards that have been put in place for responsible release of data or models that have a high risk for misuse (e.g., pre-trained language models, image generators, or scraped datasets)?
    \item[] Answer: \answerNA{}
    \item[] Justification: The submission does not release a new high-capacity foundation model or a large scraped dataset; it describes a harness-level method and provides an anonymized implementation link without claiming a dual-use model release that requires special gating beyond standard review.
    \item[] Guidelines:
    \begin{itemize}
        \item The answer \answerNA{} means that the paper poses no such risks.
        \item Released models that have a high risk for misuse or dual-use should be released with necessary safeguards to allow for controlled use of the model, for example by requiring that users adhere to usage guidelines or restrictions to access the model or implementing safety filters. 
        \item Datasets that have been scraped from the Internet could pose safety risks. The authors should describe how they avoided releasing unsafe images.
        \item We recognize that providing effective safeguards is challenging, and many papers do not require this, but we encourage authors to take this into account and make a best faith effort.
    \end{itemize}

\item {\bf Licenses for existing assets}
    \item[] Question: Are the creators or original owners of assets (e.g., code, data, models), used in the paper, properly credited and are the license and terms of use explicitly mentioned and properly respected?
    \item[] Answer: \answerYes{}
    \item[] Justification: Benchmarks and baselines are cited in the main text and appendix~\ref{app:dataset}, which also points to consolidated license and terms-of-use documentation in \url{https://anonymous.4open.science/r/ExpG-F72B}.
    \item[] Guidelines:
    \begin{itemize}
        \item The answer \answerNA{} means that the paper does not use existing assets.
        \item The authors should cite the original paper that produced the code package or dataset.
        \item The authors should state which version of the asset is used and, if possible, include a URL.
        \item The name of the license (e.g., CC-BY 4.0) should be included for each asset.
        \item For scraped data from a particular source (e.g., website), the copyright and terms of service of that source should be provided.
        \item If assets are released, the license, copyright information, and terms of use in the package should be provided. For popular datasets, \url{paperswithcode.com/datasets} has curated licenses for some datasets. Their licensing guide can help determine the license of a dataset.
        \item For existing datasets that are re-packaged, both the original license and the license of the derived asset (if it has changed) should be provided.
        \item If this information is not available online, the authors are encouraged to reach out to the asset's creators.
    \end{itemize}

\item {\bf New assets}
    \item[] Question: Are new assets introduced in the paper well documented and is the documentation provided alongside the assets?
    \item[] Answer: \answerYes{}
    \item[] Justification: The anonymized repository \url{https://anonymous.4open.science/r/ExpG-F72B} ships code and experiment artifacts with README documentation; the method, prompts (appendix~\ref{prompt}), and implementation notes (appendix~\ref{app:tried}) further describe the pipeline.
    \item[] Guidelines:
    \begin{itemize}
        \item The answer \answerNA{} means that the paper does not release new assets.
        \item Researchers should communicate the details of the dataset\slash code\slash model as part of their submissions via structured templates. This includes details about training, license, limitations, etc. 
        \item The paper should discuss whether and how consent was obtained from people whose asset is used.
        \item At submission time, remember to anonymize your assets (if applicable). You can either create an anonymized URL or include an anonymized zip file.
    \end{itemize}

\item {\bf Crowdsourcing and research with human subjects}
    \item[] Question: For crowdsourcing experiments and research with human subjects, does the paper include the full text of instructions given to participants and screenshots, if applicable, as well as details about compensation (if any)? 
    \item[] Answer: \answerNA{}
    \item[] Justification: The study uses automated benchmarks and LLM judges; it does not involve crowdsourcing or research with human subjects.
    \item[] Guidelines:
    \begin{itemize}
        \item The answer \answerNA{} means that the paper does not involve crowdsourcing nor research with human subjects.
        \item Including this information in the supplemental material is fine, but if the main contribution of the paper involves human subjects, then as much detail as possible should be included in the main paper. 
        \item According to the NeurIPS Code of Ethics, workers involved in data collection, curation, or other labor should be paid at least the minimum wage in the country of the data collector. 
    \end{itemize}

\item {\bf Institutional review board (IRB) approvals or equivalent for research with human subjects}
    \item[] Question: Does the paper describe potential risks incurred by study participants, whether such risks were disclosed to the subjects, and whether Institutional Review Board (IRB) approvals (or an equivalent approval/review based on the requirements of your country or institution) were obtained?
    \item[] Answer: \answerNA{}
    \item[] Justification: The research does not include human subjects or crowdsourcing.
    \item[] Guidelines:
    \begin{itemize}
        \item The answer \answerNA{} means that the paper does not involve crowdsourcing nor research with human subjects.
        \item Depending on the country in which research is conducted, IRB approval (or equivalent) may be required for any human subjects research. If you obtained IRB approval, you should clearly state this in the paper. 
        \item We recognize that the procedures for this may vary significantly between institutions and locations, and we expect authors to adhere to the NeurIPS Code of Ethics and the guidelines for their institution. 
        \item For initial submissions, do not include any information that would break anonymity (if applicable), such as the institution conducting the review.
    \end{itemize}

\item {\bf Declaration of LLM usage}
    \item[] Question: Does the paper describe the usage of LLMs if it is an important, original, or non-standard component of the core methods in this research? Note that if the LLM is used only for writing, editing, or formatting purposes and does \emph{not} impact the core methodology, scientific rigor, or originality of the research, declaration is not required.
    %this research? 
    \item[] Answer: \answerYes{}
    \item[] Justification: $\mathrm{LLM}_{\mathrm{evaluator}}$ and $\mathrm{LLM}_{\mathrm{summarizer}}$ (LLM-as-judge and summarization) are core to experience acquisition and distillation; section~\ref{method} and appendix~\ref{prompt} specify their roles and prompt templates.
    \item[] Guidelines:
    \begin{itemize}
        \item The answer \answerNA{} means that the core method development in this research does not involve LLMs as any important, original, or non-standard components.
        \item Please refer to our LLM policy in the NeurIPS handbook for what should or should not be described.
    \end{itemize}

\end{enumerate}



\end{document}